% Capítulo: La fidelidad de Dios (generado desde 15_Fiel/15_Fiel_capitulo.tex)
% Texto ancla: Josué 21:43–45 y 23:14. Nombre: ha'El hane'emán (Dt 7:9). Clasificación: Relación; incomunicable en su forma absoluta, comunicable como reflejo.

\section{Introducción: una palabra que no cae}

La fidelidad escasea entre los hombres, y precisamente por eso resplandece en Dios. Pink traza un diagnóstico severo: ``La infidelidad es uno de los pecados más sobresalientes de estos días malos''. En los negocios ya no se honra la palabra empeñada; los votos matrimoniales se rompen ``con tan poca consideración como cuando se descarta un vestido viejo'', y ni el lector ni el escritor ``pueden pretender inmunidad completa de este terrible pecado''. De ahí que resulte ``indescriptiblemente bendito'' levantar los ojos y ``contemplar a Aquel que es fiel'' (Pink, 2020, cap. 10, p. 53).

La definición de esta serie cabe en una línea: \destacar{Dios cumple todo lo que promete; nunca falla}. La Escritura la expresa con una imagen física. Una promesa es una palabra lanzada hacia el futuro, y si no se cumple, ``cae'' al suelo. De la palabra de Dios se afirma, una y otra vez, justamente lo contrario: que no cayó.

\hebverso{לֹא־נָפַל דָּבָר מִכֹּל הַדָּבָר הַטּוֹב אֲשֶׁר־דִּבֶּר יְהוָה אֶל־בֵּית יִשְׂרָאֵל הַכֹּל בָּא}
\begin{verso}{Josué 21:45}
No faltó palabra de todas las buenas promesas que \SENOR{} había hecho a la casa de Israel; todo se cumplió.
\end{verso}

El verso que une esta obra invita a conocer a Dios: ``alábese en esto el que se hubiere de alabar: en entenderme y conocerme'' (Jer 9:24 [9:23]). Deuteronomio emplea el mismo verbo, \heb{יָדַע} (\emph{yadá}), al presentar el Nombre de este atributo: ``\emph{Conoce}, pues, que \SENOR{} tu Dios es Dios, Dios fiel'' (Dt 7:9). En la Biblia, conocer a Dios incluye saber que su palabra no cae.

El estudio avanza en cuatro pasos. Primero escucha lo que dice la Escritura cuando llama fiel a Dios (el texto ancla, la palabra y el Nombre); luego, cómo lo canta y lo narra Israel (los Salmos, Salomón y la tradición judía); después, cómo lo lee Pablo a la luz de Cristo; y, por último, qué pide al creyente.

%===================================================================
\section{Texto ancla: Josué 21:43--45 y 23:14}\label{sec:josue-fiel}

\subsection{El final de una larga espera}
La promesa de la tierra recorre todo el Pentateuco. A Abram, recién llegado a Canaán, Dios le dijo: ``A tu descendencia daré esta tierra'' (Gn 12:7). Más tarde, en el pacto entre los animales partidos, le anunció también que la promesa tendría que esperar:

\begin{verso}{Génesis 15:13, 16}
Entonces \SENOR{} dijo a Abram: Ten por cierto que tu descendencia morará en tierra ajena, y será esclava allí, y será oprimida cuatrocientos años\ldots{} Y en la cuarta generación volverán acá; porque aún no ha llegado a su colmo la maldad del amorreo hasta aquí.
\end{verso}

Pink toma este texto como prueba: los descendientes de Abraham ``gimieron en Egipto entre los hornos de ladrillos. ¿Había olvidado Dios su promesa? ¡De hecho, no!'' (Pink, 2020, cap. 10, p. 54). El Éxodo registra el cumplimiento con exactitud de calendario: ``pasados los cuatrocientos treinta años, en el mismo día todas las huestes de \SENOR{} salieron de la tierra de Egipto'' (Éx 12:41). Más de cuarenta años después, el libro de Josué cierra el reparto de la tierra con un balance:

\begin{verso}{Josué 21:43–45}
De esta manera dio \SENOR{} a Israel toda la tierra que había jurado dar a sus padres, y la poseyeron y habitaron en ella. Y \SENOR{} les dio reposo alrededor, conforme a todo lo que había jurado a sus padres; y ninguno de todos sus enemigos pudo hacerles frente, porque \SENOR{} entregó en sus manos a todos sus enemigos. No faltó palabra de todas las buenas promesas que \SENOR{} había hecho a la casa de Israel; todo se cumplió.
\end{verso}

Los tres versos siguen un mismo movimiento: lo que Dios ``había jurado'' (vv. 43, 44) y lo que había ``hecho'', literalmente ``hablado'' (v. 45), se convierte en hechos concretos: ``dio'', ``poseyeron'', ``habitaron'', ``reposo'', ``entregó''. Conviene notar a quién atribuye el narrador la fidelidad. No dice que Israel fue fiel; dice que lo fue Dios.

\begin{filologia}{Jos 21:45}
\heb{לֹא־נָפַל דָּבָר} \emph{lo nafal davar} & \heb{נפל} & ``no cayó una palabra''. La RVR1960 traduce ``no faltó''. El verbo \emph{nafal} es ``caer'': la palabra que no se cumple cae al suelo como una flecha que no llega. \\
\heb{הַדָּבָר הַטּוֹב} \emph{hadavar hatov} & \heb{טוב} & ``la palabra buena'', en singular colectivo: la RVR1960 traduce ``las buenas promesas''. El hebreo bíblico no tiene un sustantivo propio para ``promesa''; dice ``palabra''. \\
\heb{הַכֹּל בָּא} \emph{hakol ba} & \heb{בוא} & ``todo vino'', ``todo llegó'': la RVR1960 traduce ``todo se cumplió''. La promesa es una palabra que viaja hacia el futuro y llega. \\
\grc{οὐ διέπεσεν} & \grc{διαπίπτω} & la Septuaginta conserva la imagen: ``no se cayó'' ninguna de las palabras buenas. \\
\end{filologia}

\precision{el libro de Josué afirma a la vez que Dios dio ``toda la tierra'' (21:43) y que quedaba ``aún mucha tierra por poseer'' (13:1; cf. Jue 1). Los comentarios resuelven la tensión distinguiendo dos planos: Dios cumplió lo que había prometido (la tierra entregada, el reposo, los enemigos sometidos), mientras que la ocupación completa dependía de la obediencia de las tribus, que no echaron a todos los cananeos (MacArthur, 2011; Henry, 1999). Hebreos retoma la misma tensión desde otro ángulo: ``si Josué les hubiera dado el reposo, no hablaría después de otro día. Por tanto, queda un reposo para el pueblo de Dios'' (He 4:8--9).}

\subsection{“Reconoced con todo vuestro corazón”}
Ya anciano, Josué reúne a Israel y repite la frase en primera persona, como quien deja un testamento:

\hebverso{כִּי לֹא־נָפַל דָּבָר אֶחָד מִכֹּל הַדְּבָרִים הַטּוֹבִים אֲשֶׁר דִּבֶּר יְהוָה אֱלֹהֵיכֶם עֲלֵיכֶם \\ הַכֹּל בָּאוּ לָכֶם לֹא־נָפַל מִמֶּנּוּ דָּבָר אֶחָד}
\begin{verso}{Josué 23:14}
Y he aquí que yo estoy para entrar hoy por el camino de toda la tierra; reconoced, pues, con todo vuestro corazón y con toda vuestra alma, que no ha faltado una palabra de todas las buenas palabras que \SENOR{} vuestro Dios había dicho de vosotros; todas os han acontecido, no ha faltado ninguna de ellas.
\end{verso}

La misma frase, \emph{lo nafal davar ejad}, ``no cayó una sola palabra'', abre y cierra el verso, y entre ambas queda el centro positivo: ``todas os han acontecido''. Este pequeño quiasmo encierra la experiencia de toda una generación. Por eso Josué no pide a Israel que crea algo nuevo, sino que \emph{reconozca} (\heb{וִידַעְתֶּם}, otra vez el verbo \emph{yadá}) lo que ya vio con sus propios ojos.

\subsection{Fiel también en la advertencia}
El verso siguiente impide reducir la fidelidad a un simple optimismo:

\begin{verso}{Josué 23:15}
Pero así como ha venido sobre vosotros toda palabra buena que \SENOR{} vuestro Dios os había dicho, también traerá \SENOR{} sobre vosotros toda palabra mala, hasta destruiros de sobre la buena tierra que \SENOR{} vuestro Dios os ha dado.
\end{verso}

``Toda palabra buena'' y ``toda palabra mala'' proceden del mismo Dios, que cumple las dos. Pink lo formula así: la Palabra ``no sólo abunda en ilustraciones de su fidelidad en el cumplimiento de sus promesas, sino que también registra numerosos ejemplos de su fidelidad en el cumplimiento de sus amenazas'' (Pink, 2020, cap. 10, p. 56). La historia posterior de Israel, hasta el destierro, confirma ambas mitades de Jos 23:14--15.

%===================================================================
\section{La palabra: lo que no falla}

\subsection{\emph{Ne'emán}: firme, seguro, fiel}
El adjetivo hebreo para ``fiel'' es \heb{נֶאֱמָן} (\emph{ne'emán}), participio nifal de la raíz \heb{א־מ־ן} (\emph{amán}), la misma de \emph{emet} (verdad) y de ``amén''. El capítulo sobre la verdad estudió esa raíz en sus usos concretos, que son sostener, cargar y ser firme (véase el capítulo 13). \emph{Ne'emán} designa lo que ha quedado firme: lo que aguanta y dura, aquello en que se puede confiar. Los usos concretos del término lo dejan ver con claridad:

\begin{center}\small\renewcommand{\arraystretch}{1.18}
\begin{tabular}{@{}L{3.6cm}L{7.2cm}L{3.4cm}@{}}
\toprule
\textbf{Forma} & \textbf{Uso} & \textbf{Texto} \\ \midrule
\heb{בְּמָקוֹם נֶאֱמָן} & ``en lugar firme'': la estaca clavada que no se suelta & Is 22:23 \\
\heb{מֵימָיו נֶאֱמָנִים} & ``sus aguas serán seguras'': el agua que no se corta & Is 33:16 \\
\heb{נֶאֱמָנוֹת} & plagas ``permanentes'', que no se van & Dt 28:59 \\
\heb{וְעֵד בַּשַּׁחַק נֶאֱמָן} & ``testigo fiel en el cielo'': la luna, que vuelve cada mes & Sal 89:37 [89:38] \\
\heb{נֶאֱמָן הוּא} & Moisés, ``fiel en toda mi casa'' & Nm 12:7 \\
\heb{נֶאֱמָן לְנָבִיא} & Samuel, ``fiel profeta'' & 1~S 3:20 \\
\heb{לְבָבוֹ נֶאֱמָן} & el corazón de Abraham, hallado ``fiel'' & Neh 9:8 \\
\heb{קִרְיָה נֶאֱמָנָה} & Jerusalén, ``ciudad fiel'' que se volvió ramera & Is 1:21 \\
\heb{הָאֵל הַנֶּאֱמָן} & ``el Dios fiel'' & Dt 7:9 \\
\bottomrule
\end{tabular}\par
{\footnotesize\color{gris}\textit{Tabla 1.} Usos de \emph{ne'emán}. \textit{Nota.} Textos de la RVR1960 y la BHS; numeración española [hebrea].}
\end{center}

La estaca, el agua y la luna dicen lo mismo que el profeta y el siervo: \emph{ne'emán} es aquello que no falla cuando uno se apoya en ello. Aplicado a Dios, entonces, el adjetivo no describe un sentimiento, sino una conducta comprobada a lo largo del tiempo.

\subsection{\emph{Emuná}: la fidelidad que dura}
El sustantivo \heb{אֱמוּנָה} (\emph{emuná}) aparece unas 49 veces en el Antiguo Testamento, y su primer uso es físico: las manos de Moisés, sostenidas por Aarón y Hur, tuvieron ``firmeza hasta que se puso el sol'' (Éx 17:12). Cuando se aplica a Dios, en cambio, la \emph{emuná} se mide en generaciones y en cielos:

\hebverso{יְהוָה בְּהַשָּׁמַיִם חַסְדֶּךָ אֱמוּנָתְךָ עַד־שְׁחָקִים}
\begin{verso}{Salmo 36:5 [36:6]}
\SENOR{}, hasta los cielos llega tu misericordia, Y tu fidelidad alcanza hasta las nubes.
\end{verso}

\begin{verso}{Salmo 119:90}
De generación en generación es tu fidelidad; Tú afirmaste la tierra, y subsiste.
\end{verso}

El Salmo 119:90 enlaza la fidelidad de Dios con la estabilidad del mundo: la tierra ``subsiste'' (\heb{וַתַּעֲמֹד}, ``se mantiene en pie'') porque Dios es fiel. Se trata de la misma promesa que Dios dio después del diluvio y que Pink presenta como una prueba renovada cada año: ``Mientras la tierra permanezca, no cesarán la sementera y la siega, el frío y el calor, el verano y el invierno, y el día y la noche'' (Gn 8:22; Pink, 2020, cap. 10, p. 54). Jeremías invocará ese orden del día y de la noche como garantía del pacto con David (Jer 33:20--21).

La \emph{emuná} es también una vestidura. Isaías describe al Mesías ceñido de ella: ``Y será la justicia cinto de sus lomos, y la fidelidad ceñidor de su cintura'' (Is 11:5); y el salmista contempla a Dios envuelto en ella: ``Poderoso eres, \SENOR, y tu fidelidad te rodea'' (Sal 89:8 [89:9]). Pink lo comenta con una imagen: ``Es como si Dios estuviera revestido de ella'' (Pink, 2020, cap. 10, p. 53).

\subsection{Los pares fijos: \emph{jésed} y \emph{emuná}}
En la poesía hebrea la fidelidad casi nunca aparece sola. Forma pareja con \heb{חֶסֶד} (\emph{jésed}), el amor leal del pacto que el capítulo 14 estudió como misericordia. Dt 7:9 ya los une en el Nombre, pues Dios ``guarda el pacto y la misericordia'', y los Salmos los distribuyen entre los dos miembros del paralelismo:

\hebverso{לְהַגִּיד בַּבֹּקֶר חַסְדֶּךָ וֶאֱמוּנָתְךָ בַּלֵּילוֹת}
\begin{verso}{Salmo 92:2 [92:3]}
Anunciar por la mañana tu misericordia, Y tu fidelidad cada noche.
\end{verso}

``Mañana'' y ``noches'' forman un merismo, en el que los dos extremos nombran todo el tiempo, del mismo modo que ``cielos'' y ``nubes'' nombran todo el espacio en Sal 36:5. La \emph{jésed} es el amor que se compromete; la \emph{emuná}, ese mismo amor sostenido en el tiempo. Así lo cantan otros salmos: ``Porque \SENOR{} es bueno; para siempre es su misericordia, y su verdad por todas las generaciones'' (Sal 100:5); ``y la fidelidad de \SENOR{} es para siempre'' (Sal 117:2).

\precision{en Sal 100:5 y 117:2 el hebreo dice \emph{emuná} y \emph{emet}, respectivamente; la RVR1960 traduce la primera ``verdad'' y la segunda ``fidelidad''. Como se vio en el capítulo 13, ambas palabras proceden de la misma raíz, y las versiones eligen uno u otro matiz según el contexto.}

\subsection{El espectro completo}
En la RVR1960, ``fiel'', ``fieles'', ``fielmente'', ``fidelidad'', ``infiel'', ``infieles'' e ``infidelidad'' aparecen 122 veces en 118 versículos (60 en el Antiguo Testamento y 62 en el Nuevo), sin contar los títulos de sección. Tras esas palabras castellanas se esconden varios términos del original:

\begin{center}\small\renewcommand{\arraystretch}{1.16}
\begin{tabular}{@{}L{3.6cm}L{1.3cm}L{5.2cm}L{4.4cm}@{}}
\toprule
\textbf{Término} & \textbf{Veces} & \textbf{Traducción en la RVR1960} & \textbf{Textos} \\ \midrule
\heb{נֶאֱמָן} \emph{ne'emán} & $\approx$\,34 & ``fiel'' (la mayoría), ``firme'', ``seguro'', ``permanente'', ``duradero'' & Dt 7:9; Is 49:7; Sal 89:37 \\
\heb{אֱמוּנָה} \emph{emuná} & 49 & ``fidelidad'' (13), ``verdad'' (18), ``fielmente'', ``firmeza'', ``fe'' & Sal 36:5; Lam 3:23; Hab 2:4 \\
\heb{אֱמֶת} \emph{emet} & 127 & ``verdad''; ``fidelidad'' en 2~S 15:20; 2~Cr 32:1; Sal 117:2; 138:2 & Sal 117:2 \\
\heb{אֱמוּנִים} \emph{emunim} & 5 & ``fieles'' (2~S 20:19; Sal 12:1; 31:23), ``fiel'' (Pr 13:17) & Sal 31:23 \\
\heb{אֵמֻן} \emph{emún} & 1 & ``infieles'' (``hijos sin fidelidad'') & Dt 32:20 \\
\heb{אֲמָנָה} \emph{amaná} & 2 & ``fiel promesa'' (pacto escrito) & Neh 9:38 [10:1] \\
arameo \heb{מְהֵימַן} \emph{mehemán} & 2 & ``fiel'' (la interpretación; Daniel) & Dn 2:45; 6:4 \\
\heb{חֶסֶד} \emph{jésed} & --- & ``la fidelidad de tu juventud'' & Jer 2:2 \\
\heb{מִשְׁמֶרֶת} \emph{mishméret} & --- & ``guardaban fidelidad'' (lo que se custodia) & 1~Cr 12:29 [12:30]; Ez 48:11 \\
\heb{מִלֵּא אַחֲרֵי} \emph{milé ajaré} & --- & ``ha seguido fielmente'' (``llenó tras'') & Dt 1:36 \\ \midrule
\grc{πιστός} \emph{pistós} & 67 & ``fiel'' ($\approx$\,55), ``creyente'', ``digno de crédito'' & 1~Co 1:9; 1~Ts 5:24; Ap 19:11 \\
\grc{πίστις} \emph{pístis} & 243 & ``fe''; ``fidelidad'' en Ro 3:3; ``fieles'' en Tit 2:10 & Ro 3:3; Gá 5:22 \\
\grc{ἀπιστέω}, \grc{ἄπιστος}, \grc{ἀπιστία} & 8; 23; 11 & ``ser incrédulo'', ``infiel'', ``incredulidad'' & Ro 3:3; 2~Ti 2:13 \\ \midrule
\heb{מַעַל} \emph{ma'al} & --- & ``infidelidad'', ``ser infiel'' (traición a lo sagrado) & Nm 5:12, 27; 2~Cr 29:19 \\
\heb{בָּגַד} \emph{bagad} & --- & ``la esposa infiel'', ``prevaricar'', ``traicionar'' & Jer 3:20; Job 6:15 \\
\heb{כָּזָב}, \heb{אַכְזָב} \emph{kazav}, \emph{ajzav} & --- & ``mentira''; ``cosa ilusoria'', el arroyo que se seca & Jer 15:18; Mi 1:14 \\
\bottomrule
\end{tabular}\par
{\footnotesize\color{gris}\textit{Tabla 2.} El espectro de ``fiel'' y ``fidelidad''. \textit{Nota.} Conteo propio con la BHS y NA27 con números Strong, y la RVR1960 sin títulos de sección. \emph{Ne'emán}: solo el participio nifal.}
\end{center}

\precision{los conteos son propios y conviene verificarlos en una concordancia (e-Sword). La cifra de \emph{emuná} (49) coincide con la del capítulo 13. Para \emph{jésed}, \emph{mishméret} y los contrarios no se da cifra, porque solo algunos de sus usos se traducen con la familia de ``fiel''.}

Del espectro se desprenden tres enseñanzas. La primera es que, en hebreo, fidelidad y verdad pertenecen a la misma familia: lo fiel es lo que, con el tiempo, resulta verdadero. La segunda, que en griego \grc{πιστός} abarca las dos caras de la relación, el que es digno de confianza (Dios es \emph{pistós}) y el que confía (el creyente es \emph{pistós}). La tercera, que lo contrario de la fidelidad no es ante todo un error, sino una traición: \emph{ma'al} y \emph{bagad} son las palabras del adulterio y del sacrilegio.

\subsection{Los verbos de cumplir: no caer, llenar, levantar}
El hebreo no dispone de un único verbo para ``cumplir una promesa''; recurre, más bien, a imágenes:

\begin{center}\small\renewcommand{\arraystretch}{1.2}
\begin{tabular}{@{}L{3.6cm}L{6.6cm}L{4.2cm}@{}}
\toprule
\textbf{Expresión} & \textbf{Imagen} & \textbf{Textos} \\ \midrule
\heb{לֹא נָפַל דָּבָר} & la palabra ``no cayó'' al suelo & Jos 21:45; 23:14; 1~R 8:56; 2~R 10:10 \\
\heb{לֹא הִפִּיל אָרְצָה} & ``no dejó caer a tierra'' ninguna palabra & 1~S 3:19 \\
\heb{מִלֵּא} (piel) & ``llenó'': lo dicho con la boca, la mano lo llena & 1~R 8:15, 24 \\
\heb{הֵקִים} (hifil de \emph{qum}) & ``levantó'', puso en pie la palabra & Dt 9:5; 1~R 8:20; Neh 9:8; Jer 33:14 \\
\heb{יֵאָמֵן} (nifal de \emph{amán}) & ``sea afirmada'' la palabra & 1~R 8:26 \\
\heb{שֹׁקֵד עַל־דְּבָרִי} & ``vigilo mi palabra'' para cumplirla & Jer 1:12 \\
\heb{לֹא יָשׁוּב רֵיקָם} & la palabra ``no volverá vacía'' & Is 55:11 \\
\bottomrule
\end{tabular}\par
{\footnotesize\color{gris}\textit{Tabla 3.} Cómo dice el hebreo ``cumplir''. \textit{Nota.} Elaboración propia con la BHS.}
\end{center}

Estas imágenes forman un conjunto coherente. La palabra que no se cumple \emph{cae}; la que se cumple se \emph{levanta} y queda en pie. Lo que la boca dice es como un recipiente vacío hasta que la mano lo \emph{llena}, y el Nuevo Testamento hereda precisamente esta última imagen en su fórmula de cumplimiento: ``para que se cumpliese'' (\grc{ἵνα πληρωθῇ}), es decir, ``para que se llenase''.

A Jeremías se le comunica el mismo mensaje mediante un juego de palabras:

\begin{verso}{Jeremías 1:11–12}
La palabra de \SENOR{} vino a mí, diciendo: ¿Qué ves tú, Jeremías? Y dije: Veo una vara de almendro. Y me dijo \SENOR: Bien has visto; porque yo apresuro mi palabra para ponerla por obra.
\end{verso}

\begin{filologia}{Jer 1:11--12}
\heb{מַקֵּל שָׁקֵד} \emph{maqel shaqed} & \heb{שׁקד} & ``vara de almendro''. El almendro es el primer árbol que florece al final del invierno: ``el que vela''. \\
\heb{שֹׁקֵד אֲנִי} \emph{shoqed aní} & \heb{שׁקד} & ``yo estoy velando'', vigilante. La RVR1960 traduce ``apresuro''. \\
\heb{עַל־דְּבָרִי לַעֲשֹׂתוֹ} & \heb{עשׂה} & ``sobre mi palabra, para hacerla''. \\
\end{filologia}

\emph{Shaqed} (almendro) y \emph{shoqed} (vigilante) suenan casi igual, y el parecido encierra la enseñanza. Dios no pronuncia una palabra para luego abandonarla: la vigila hasta que se cumple, como el almendro que, al florecer antes que los demás árboles, anuncia que el invierno ha terminado.

\subsection{Del hebreo a Pablo}
La Septuaginta traduce \emph{ne'emán} con \grc{πιστός} y forja así la fórmula que Pablo repetirá. En Dt 7:9 escribe \grc{θεὸς πιστός}, ``Dios fiel''; en Dt 32:4, donde el hebreo dice \emph{El emuná}, vuelve a escribir \grc{θεὸς πιστός, καὶ οὐκ ἔστιν ἀδικία}; y en Is 49:7 traduce \grc{πιστός ἐστιν ὁ ἅγιος Ἰσραηλ}, ``fiel es el Santo de Israel''. De ahí procede la confesión de Pablo: \grc{πιστὸς ὁ θεός}, ``Fiel es Dios'' (1~Co 1:9; 10:13).

\begin{center}
\begin{tikzpicture}[font=\small,node distance=0.9cm]
 \node[draw=azul,line width=0.8pt,fill=papiro,rounded corners=3pt,text width=3.6cm,align=center,inner sep=6pt] (h) {\textbf{Hebreo}\\ \heb{הָאֵל הַנֶּאֱמָן} (Dt 7:9)\\ \heb{אֵל אֱמוּנָה} (Dt 32:4)\\ \heb{נֶאֱמָן קְדֹשׁ יִשְׂרָאֵל} (Is 49:7)};
 \node[draw=azul,line width=0.8pt,fill=papiro,rounded corners=3pt,text width=3.6cm,align=center,inner sep=6pt,right=of h] (l) {\textbf{Septuaginta}\\ \grc{θεὸς πιστός}\\ \grc{θεὸς πιστός}\\ \grc{πιστός ἐστιν ὁ ἅγιος}};
 \node[draw=dorado,line width=1.6pt,fill=azul,text=papiro,rounded corners=3pt,text width=3.6cm,align=center,inner sep=6pt,right=of l] (p) {\textbf{Pablo}\\ \grc{πιστὸς ὁ θεός}\\ (1 Co 1:9; 10:13; 2 Co 1:18)\\ \grc{πιστὸς ὁ καλῶν} (1 Ts 5:24)};
 \draw[-{Latex[length=3mm]},dorado,line width=1.4pt] (h)--(l);
 \draw[-{Latex[length=3mm]},dorado,line width=1.4pt] (l)--(p);
\end{tikzpicture}\par
{\footnotesize\color{gris}\textit{Figura 1.} Del hebreo a Pablo. \textit{Nota.} Elaboración propia con la BHS, la LXX (Rahlfs-Hanhart) y NA28.}
\end{center}

En el griego clásico, \grc{πιστός} se dice del amigo o del testigo digno de crédito. En la Septuaginta y en Pablo, sin embargo, la palabra carga con el peso del hebreo: ya no significa solo ``creíble'', sino firme como la estaca y seguro como el agua que no se corta.

%===================================================================
\section{El Nombre que revela la fidelidad}\label{sec:nombre-fiel}

\subsection{\emph{Ha'El hane'emán}: “el Dios fiel”}
El Nombre de este atributo aparece en el discurso de Moisés sobre la elección de Israel. Dios no eligió a Israel por su tamaño: ``erais el más insignificante de todos los pueblos; sino por cuanto \SENOR{} os amó, y quiso guardar el juramento que juró a vuestros padres'' (Dt 7:7--8). El amor explica la elección; el juramento explica su permanencia. Sobre ese fundamento, Moisés pronuncia el Nombre:

\hebverso{וְיָדַעְתָּ כִּי־יְהוָה אֱלֹהֶיךָ הוּא הָאֱלֹהִים הָאֵל הַנֶּאֱמָן \\ שֹׁמֵר הַבְּרִית וְהַחֶסֶד לְאֹהֲבָיו וּלְשֹׁמְרֵי מִצְוֹתָו לְאֶלֶף דּוֹר}
\begin{verso}{Deuteronomio 7:9}
Conoce, pues, que \SENOR{} tu Dios es Dios, Dios fiel, que guarda el pacto y la misericordia a los que le aman y guardan sus mandamientos, hasta mil generaciones.
\end{verso}

\begin{filologia}{Dt 7:9}
\heb{וְיָדַעְתָּ} \emph{veyadatá} & \heb{ידע} & ``y conocerás'', ``conoce'': el verbo de Jer 9:24. La fidelidad es algo que se conoce. \\
\heb{הָאֵל הַנֶּאֱמָן} \emph{ha'El hane'emán} & \heb{אמן} & ``el Dios, el fiel'': con artículo, como título. LXX \grc{θεὸς πιστός}. \\
\heb{שֹׁמֵר הַבְּרִית וְהַחֶסֶד} & \heb{שׁמר} & ``que guarda el pacto y la \emph{jésed}''. El mismo verbo describe a los que ``guardan'' sus mandamientos: Dios guarda primero. \\
\heb{לְאֶלֶף דּוֹר} \emph{le'élef dor} & \heb{דור} & ``por mil generaciones'': explica los ``millares'' de Éx 20:6 y 34:7 como generaciones. \\
\end{filologia}

La expresión ``Dios fiel'' ocupa el centro del discurso. La precede la elección por amor (7:6--8) y la sigue la advertencia de que Dios ``da el pago en persona al que le aborrece'' (7:10). Como ocurría en Jos 23:14--15, la fidelidad que revela el Nombre abarca tanto la promesa como la advertencia.

\subsection{“Fiel es el Santo de Israel”}
En el segundo cántico del Siervo, Dios se dirige al ``menospreciado de alma, al abominado de las naciones'' y le anuncia que los reyes se levantarán ante él:

\hebverso{לְמַעַן יְהוָה אֲשֶׁר נֶאֱמָן קְדֹשׁ יִשְׂרָאֵל וַיִּבְחָרֶךָּ}
\begin{verso}{Isaías 49:7}
Así ha dicho \SENOR, Redentor de Israel, el Santo suyo, al menospreciado de alma, al abominado de las naciones, al siervo de los tiranos: Verán reyes, y se levantarán príncipes, y adorarán por \SENOR{}; porque fiel es el Santo de Israel, el cual te escogió.
\end{verso}

La exaltación del Siervo no tiene su causa en él, sino en Dios: \emph{ne'emán qedosh Israel}, ``fiel es el Santo de Israel''. Aquí el título de santidad (capítulo 4) se une al de fidelidad, y la unión tiene una consecuencia: el que eligió no abandona a su elegido, aunque las naciones lo desprecien. El versículo siguiente es el que Pablo aplicará a su propio ministerio: ``En tiempo aceptable te oí, y en el día de salvación te ayudé'' (Is 49:8; 2~Co 6:2).

\subsection{Otros nombres y títulos}
En el cántico de Dt 32:4, Moisés llama a Dios \heb{אֵל אֱמוּנָה} (\emph{El emuná}), título que el capítulo 13 estudió como ``Dios de verdad'' y que la Septuaginta traduce \grc{θεὸς πιστός}, ``Dios fiel''. En el Nuevo Testamento el título pasa a Cristo: el jinete del caballo blanco ``se llamaba Fiel y Verdadero'' (Ap 19:11), y el Resucitado se presenta como ``el Amén, el testigo fiel y verdadero'' (Ap 3:14).

\begin{center}\small\renewcommand{\arraystretch}{1.22}
\begin{tabular}{@{}L{3.2cm}L{3.6cm}L{2.6cm}L{5.0cm}@{}}
\toprule
\textbf{Nombre o título} & \textbf{Hebreo / griego} & \textbf{Texto} & \textbf{Qué revela} \\ \midrule
Dios fiel & \heb{הָאֵל הַנֶּאֱמָן} & Dt 7:9 & guarda pacto y \emph{jésed} por mil generaciones \\
Fiel es el Santo de Israel & \heb{נֶאֱמָן קְדֹשׁ יִשְׂרָאֵל} & Is 49:7 & no abandona al Siervo que eligió \\
Dios de fidelidad & \heb{אֵל אֱמוּנָה}; LXX \grc{θεὸς πιστός} & Dt 32:4 & su obra es perfecta; en él no hay iniquidad \\
Fiel es Dios & \grc{πιστὸς ὁ θεός} & 1~Co 1:9; 10:13 & llama y sostiene hasta el fin \\
Fiel Creador & \grc{πιστὸς κτίστης} & 1~P 4:19 & se le encomienda el alma en el sufrimiento \\
Testigo fiel & \grc{ὁ μάρτυς ὁ πιστός} & Ap 1:5; 3:14 & Cristo, el Amén de Dios \\
Fiel y Verdadero & \grc{πιστὸς καὶ ἀληθινός} & Ap 19:11 & Cristo que juzga y vence \\
\bottomrule
\end{tabular}\par
{\footnotesize\color{gris}\textit{Tabla 4.} Los nombres que revelan la fidelidad. \textit{Nota.} Elaboración propia.}
\end{center}

\precision{la tradición judía resume el principio de estos títulos en el midrash \emph{Shemot Rabá} 3:6, donde Dios dice a Moisés: ``según mis obras soy llamado''. Cada nombre nace de una obra, y ``el Dios fiel'' nace de las promesas cumplidas. La fórmula \emph{El Mélej Ne'emán}, ``Dios, Rey fiel'', es posterior y se estudia en la sección sobre el judaísmo.}

%===================================================================
\section{Definición}

La definición de la serie no es otra cosa que Josué 21:45 dicho en lenguaje actual:

\begin{tcolorbox}[colback=papiro,colframe=dorado,boxrule=1.2pt,arc=3pt]
\centering\large\itshape Dios cumple todo lo que promete; nunca falla.
\end{tcolorbox}

A la luz de los textos de este capítulo, puede desarrollarse así:

\begin{tcolorbox}[colback=papiro,colframe=azul,boxrule=1pt,arc=3pt,fonttitle=\bfseries,title=Definición]
La fidelidad de Dios es la perfección eterna por la cual Dios cumple todo lo que promete y nunca falla. Guarda su pacto y su palabra de generación en generación; cumple lo que anuncia para bendecir y también lo que anuncia para corregir; y permanece fiel aun cuando su pueblo es infiel, porque no puede negarse a sí mismo (Dt 7:9; Jos 21:45; 2~Ti 2:13).
\end{tcolorbox}

\begin{center}\small\renewcommand{\arraystretch}{1.2}
\begin{tabular}{@{}L{3.6cm}L{6.4cm}L{4.4cm}@{}}
\toprule
\textbf{Rasgo} & \textbf{Qué afirma} & \textbf{Textos} \\ \midrule
Cumple lo que dice & ninguna palabra suya cae al suelo & Jos 21:45; Nm 23:19 \\
Constante & su fidelidad no se agota con los días ni con las generaciones & Sal 119:90; Lam 3:23 \\
Fiel al pacto & guarda el pacto y la \emph{jésed} por mil generaciones & Dt 7:9; Sal 89:28, 34 \\
Fiel en la advertencia & cumple también la ``palabra mala'' y la disciplina & Jos 23:15; Sal 119:75 \\
Fiel ante el infiel & la infidelidad humana no anula la suya & Ro 3:3; 2~Ti 2:13 \\
Fiel hasta el fin & termina lo que comienza & 1~Ts 5:24; Flp 1:6 \\
Cumplida en Cristo & todas las promesas son en él ``Sí'' & 2~Co 1:20; Hch 13:32--34 \\
\bottomrule
\end{tabular}\end{center}

\textbf{Clasificación.} La fidelidad pertenece a los atributos de la \emph{relación con el ser humano}, pues responde a la pregunta ``¿cómo se revela Dios?''. Se revela en la historia, porque toda promesa necesita tiempo para cumplirse; y, sin embargo, es eterna como todas sus perfecciones: ``Para siempre, oh \SENOR, permanece tu palabra en los cielos. De generación en generación es tu fidelidad'' (Sal 119:89--90). En su forma absoluta (``nunca falla'') es \emph{incomunicable}, ya que solo de Dios puede decirse que ninguna de sus palabras cayó. Como reflejo, en cambio, es \emph{comunicable}: Moisés fue ``fiel en toda mi casa'' (Nm 12:7), y la fidelidad figura entre los frutos del Espíritu (Gá 5:22).

\precision{la definición breve es de la serie. La definición desarrollada y la clasificación (comunicable como reflejo) son una propuesta que sigue el criterio empleado en los capítulos 13 y 16, y conviene confirmarlas.}

%===================================================================
\section{El arroyo engañoso y el manantial}\label{sec:arroyo-fiel}

\subsection{Un río que se seca en verano}
En Canaán abundan los \emph{wadis}, cauces que en invierno corren llenos y en verano quedan secos. Para el caminante del desierto, un \emph{wadi} equivale a una promesa: quien lo vio con agua en invierno espera encontrarla también en verano. Job, traicionado por sus amigos, recurre a esa imagen:

\begin{verso}{Job 6:15–20}
Pero mis hermanos me traicionaron como un torrente; Pasan como corrientes impetuosas Que están escondidas por la helada, Y encubiertas por la nieve; Que al tiempo del calor son deshechas, Y al calentarse, desaparecen de su lugar; Se apartan de la senda de su rumbo, Van menguando, y se pierden. Miraron los caminantes de Temán, Los caminantes de Sabá esperaron en ellas; Pero fueron avergonzados por su esperanza; Porque vinieron hasta ellas, y se hallaron confusos.
\end{verso}

``Me traicionaron'' traduce \heb{בָּגְדוּ} (\emph{bagdú}), el verbo que Jeremías aplica a la esposa infiel (Jer 3:20). El amigo infiel se parece a un arroyo que corre lleno cuando nadie lo necesita y está seco cuando se lo busca: las caravanas de Temán y de Sabá llegaron confiadas y solo encontraron arena.

\subsection{“¿Serás para mí como cosa ilusoria?”}
El hebreo tiene un nombre para ese arroyo: \heb{אַכְזָב} (\emph{ajzav}), de la raíz \heb{כזב}, ``mentir, defraudar''. El arroyo que se seca es, literalmente, ``un mentiroso''. Y Jeremías, en su crisis más honda, llega a dirigir esa palabra a Dios mismo:

\hebverso{הָיוֹ תִהְיֶה לִי כְּמוֹ אַכְזָב מַיִם לֹא נֶאֱמָנוּ}
\begin{verso}{Jeremías 15:18}
¿Por qué fue perpetuo mi dolor, y mi herida desahuciada no admitió curación? ¿Serás para mí como cosa ilusoria, como aguas que no son estables?
\end{verso}

\begin{filologia}{Jer 15:18}
\heb{כְּמוֹ אַכְזָב} \emph{kemó ajzav} & \heb{כזב} & ``como un engañoso'': el arroyo que defrauda. La RVR1960 traduce ``cosa ilusoria''. \\
\heb{מַיִם לֹא נֶאֱמָנוּ} & \heb{אמן} & ``aguas que no son fieles'', que no se mantienen: el mismo verbo de \emph{ne'emán}. La RVR1960 traduce ``aguas que no son estables''. \\
\grc{ὡς ὕδωρ ψευδὲς οὐκ ἔχον πίστιν} & \grc{πίστις} & la Septuaginta: ``como agua mentirosa que no tiene fidelidad''. \\
\end{filologia}

La pregunta es atrevida, y sin embargo forma parte de la Escritura. Dios no la censura, sino que responde: ``Si te convirtieres, yo te restauraré\ldots{} y te pondré en este pueblo por muro fortificado de bronce'' (Jer 15:19--20). Miqueas emplea la misma palabra en un juego con el nombre de una ciudad: ``las casas de Aczib (\heb{אַכְזִיב}) serán para engaño (\heb{לְאַכְזָב}) a los reyes de Israel'' (Mi 1:14). La ciudad cuyo nombre sonaba a ``Engaño'' terminaría engañando a los reyes que confiaran en ella.

\subsection{El manantial cuyas aguas nunca faltan}
Frente al arroyo engañoso, Isaías describe la vida de quien depende del \SENOR:

\hebverso{וְהָיִיתָ כְּגַן רָוֶה וּכְמוֹצָא מַיִם אֲשֶׁר לֹא־יְכַזְּבוּ מֵימָיו}
\begin{verso}{Isaías 58:11}
\SENOR{} te pastoreará siempre, y en las sequías saciará tu alma, y dará vigor a tus huesos; y serás como huerto de riego, y como manantial de aguas, cuyas aguas nunca faltan.
\end{verso}

Donde la RVR1960 dice ``cuyas aguas nunca faltan'', el hebreo escribe \heb{אֲשֶׁר לֹא־יְכַזְּבוּ מֵימָיו}: ``cuyas aguas no mienten'', con el mismo verbo de \emph{ajzav}. El manantial no defrauda en la sequía; su agua está ahí cuando se la necesita. En la misma línea, Is 33:16 promete al justo \heb{מֵימָיו נֶאֱמָנִים}, ``sus aguas fieles'' (RVR1960: ``seguras''), que la Septuaginta traduce \grc{τὸ ὕδωρ αὐτοῦ πιστόν}.

\begin{center}
\begin{tikzpicture}[font=\small]
 % Arroyo engañoso
 \begin{scope}
  \node[font=\bfseries,text=gris] at (2.6,3.3) {El arroyo engañoso (\emph{ajzav})};
  \fill[papiro] (0,0) rectangle (5.2,2.8);
  \draw[gris,line width=0.8pt] (0,0) rectangle (5.2,2.8);
  \node[font=\scriptsize,text=azul] at (1.3,2.45) {invierno};
  \draw[azul!70,line width=4pt,line cap=round] (0.3,1.9) .. controls (0.9,1.6) and (1.5,2.1) .. (2.3,1.8);
  \node[font=\scriptsize,text=gris] at (3.9,2.45) {verano};
  \draw[dorado!80!gris,line width=4pt,line cap=round,dash pattern=on 3pt off 3pt] (2.9,1.9) .. controls (3.5,1.6) and (4.1,2.1) .. (4.9,1.8);
  \node[text width=4.8cm,align=center,font=\scriptsize] at (2.6,0.75) {corre lleno cuando no se lo necesita;\\ seco cuando se lo busca\\ (Job 6:15--20; Jer 15:18; Mi 1:14)};
 \end{scope}
 % Manantial
 \begin{scope}[xshift=6.2cm]
  \node[font=\bfseries,text=azul] at (2.6,3.3) {El manantial fiel (\emph{ne'emán})};
  \fill[azul!8] (0,0) rectangle (5.2,2.8);
  \draw[azul,line width=1.2pt] (0,0) rectangle (5.2,2.8);
  \node[font=\scriptsize,text=azul] at (1.3,2.45) {invierno};
  \node[font=\scriptsize,text=azul] at (3.9,2.45) {verano};
  \draw[azul!70,line width=4pt,line cap=round] (0.3,1.9) .. controls (1.5,1.6) and (3.4,2.1) .. (4.9,1.8);
  \node[text width=4.8cm,align=center,font=\scriptsize] at (2.6,0.75) {``cuyas aguas no mienten'';\\ ``sus aguas serán fieles''\\ (Is 58:11; 33:16; Jer 2:13; 17:13)};
 \end{scope}
\end{tikzpicture}\par
{\footnotesize\color{gris}\textit{Figura 2.} El arroyo engañoso y el manantial fiel. \textit{Nota.} Elaboración propia.}
\end{center}

\subsection{La fuente de agua viva}
Jeremías da al \SENOR{} el nombre que corresponde a esta imagen. El pecado de Israel no consistió solo en buscar otra agua, sino en cambiar el manantial por depósitos que la pierden:

\begin{verso}{Jeremías 2:13}
Porque dos males ha hecho mi pueblo: me dejaron a mí, fuente de agua viva, y cavaron para sí cisternas, cisternas rotas que no retienen agua.
\end{verso}

El profeta vuelve sobre ello más adelante: ``Oh \SENOR, esperanza de Israel\ldots{} dejaron a \SENOR, manantial de aguas vivas'' (Jer 17:13). Y Jesús retoma la promesa cuando dice que ``el agua que yo le daré será en él una fuente de agua que salte para vida eterna'' (Jn 4:14).

La imagen condensa todo el atributo. Lo infiel no siempre es lo que nunca dio nada; con frecuencia es lo que dio en invierno y falta en verano. Preguntar por la fidelidad de Dios equivale a formular la pregunta de Jeremías: ¿será Dios como el \emph{wadi} en tiempo de sequía? La Escritura responde con el manantial, cuya agua está presente cuando se la busca.

\precision{la elección de esta imagen como eje del capítulo es una decisión de esta serie. Los textos (Job 6; Jer 15:18; Is 58:11; Mi 1:14) y el sentido de \emph{ajzav} como arroyo que se seca están en los diccionarios (BDB; HALOT).}

%===================================================================
\section{La fidelidad cantada}

\subsection{El Salmo 89: siete veces fidelidad, siete veces \emph{jésed}}
El Salmo 89 es el gran poema de la fidelidad de Dios, y se abre con un voto:

\hebverso{חַסְדֵי יְהוָה עוֹלָם אָשִׁירָה לְדֹר וָדֹר אוֹדִיעַ אֱמוּנָתְךָ בְּפִי}
\begin{verso}{Salmo 89:1–2 [89:2–3]}
Las misericordias de \SENOR{} cantaré perpetuamente; De generación en generación haré notoria tu fidelidad con mi boca. Porque dije: Para siempre será edificada misericordia; En los cielos mismos afirmarás tu verdad.
\end{verso}

El salmista cumple su voto mediante una repetición que el oyente hebreo podía percibir: en el salmo, \emph{emuná} aparece siete veces y \emph{jésed} otras siete.

\begin{center}\small\renewcommand{\arraystretch}{1.18}
\begin{tabular}{@{}L{1.6cm}L{6.4cm}L{6.4cm}@{}}
\toprule
\textbf{Verso} & \textbf{\emph{emuná}} (RVR1960) & \textbf{\emph{jésed}} (RVR1960) \\ \midrule
1 [2] & ``tu fidelidad'' & ``las misericordias'' \\
2 [3] & ``tu verdad'' & ``misericordia'' \\
5 [6] & ``tu verdad'' & --- \\
8 [9] & ``tu fidelidad te rodea'' & --- \\
14 [15] & --- & ``misericordia y verdad'' (\emph{emet}) \\
24 [25] & ``mi verdad'' & ``mi misericordia'' \\
28 [29] & (\emph{ne'emán}: ``mi pacto será firme'') & ``mi misericordia'' \\
33 [34] & ``ni falsearé mi verdad'' & ``no quitaré\ldots{} mi misericordia'' \\
49 [50] & ``que juraste a David por tu verdad'' & ``tus antiguas misericordias'' \\
\bottomrule
\end{tabular}\par
{\footnotesize\color{gris}\textit{Tabla 5.} \emph{Emuná} y \emph{jésed} en el Salmo 89. \textit{Nota.} Conteo propio con la BHS; numeración española [hebrea]. La RVR1960 traduce \emph{emuná} ``fidelidad'' dos veces y ``verdad'' cinco.}
\end{center}

El dato revela cómo opera la traducción: quien lee la RVR1960 encuentra ``fidelidad'' solo dos veces, mientras que el hebreo repite la palabra siete. Además, ambas palabras aparecen juntas en el primer verso y en el último uso (v. 49), de modo que la pareja \emph{jésed} y \emph{emuná} enmarca todo el salmo.

\subsection{La promesa a David y su crisis}
En el centro del salmo está la promesa que Natán dio a David (2~S 7:12--16), ahora cantada como juramento:

\begin{verso}{Salmo 89:30–35 [89:31–36]}
Si dejaren sus hijos mi ley, Y no anduvieren en mis juicios, Si profanaren mis estatutos, Y no guardaren mis mandamientos, Entonces castigaré con vara su rebelión, Y con azotes sus iniquidades. Mas no quitaré de él mi misericordia, Ni falsearé mi verdad. No olvidaré mi pacto, Ni mudaré lo que ha salido de mis labios. Una vez he jurado por mi santidad, Y no mentiré a David.
\end{verso}

\begin{filologia}{Sal 89:33--35}
\heb{וְלֹא־אֲשַׁקֵּר בֶּאֱמוּנָתִי} & \heb{שׁקר} & ``y no mentiré en mi fidelidad'': el verbo de \emph{shéqer}, la mentira (cap. 13). RVR1960: ``ni falsearé mi verdad''. \\
\heb{וּמוֹצָא שְׂפָתַי לֹא אֲשַׁנֶּה} & \heb{שׁנה} & ``y lo que salió de mis labios no lo cambiaré''. \\
\heb{אִם־לְדָוִד אֲכַזֵּב} & \heb{כזב} & ``si a David miento'': fórmula de juramento. El verbo de \emph{ajzav}: Dios no será un arroyo engañoso para David. \\
\end{filologia}

Pink recurre a estos versos para mostrar que la disciplina no contradice la fidelidad: ``El castigo, no sólo es reconciliable con la tierna misericordia de Dios, sino que es el efecto y la expresión de la misma'' (Pink, 2020, cap. 10, p. 57). Dios promete la vara y, al mismo tiempo, promete no retirar la \emph{jésed}; ambas promesas son expresión de su fidelidad.

El salmo, sin embargo, no se detiene en el juramento. A partir del v. 38 el tono cambia bruscamente:

\begin{verso}{Salmo 89:38–39, 46, 49 [89:39–40, 47, 50]}
Mas tú desechaste y menospreciaste a tu ungido, Y te has airado con él. Rompiste el pacto de tu siervo; Has profanado su corona hasta la tierra\ldots{} ¿Hasta cuándo, oh \SENOR? ¿Te esconderás para siempre?\ldots{} Señor, ¿dónde están tus antiguas misericordias, Que juraste a David por tu verdad?
\end{verso}

\begin{center}
\begin{tikzpicture}[font=\scriptsize]
 \node[draw=azul,line width=0.9pt,fill=papiro,rounded corners=3pt,text width=3.4cm,align=center,inner sep=5pt,minimum height=1.9cm] (a) {\textbf{vv. 1--18}\\ Himno: la \emph{jésed} y la \emph{emuná} de Dios en los cielos y en la creación};
 \node[draw=azul,line width=0.9pt,fill=papiro,rounded corners=3pt,text width=3.4cm,align=center,inner sep=5pt,minimum height=1.9cm,right=0.5cm of a] (b) {\textbf{vv. 19--37}\\ Oráculo: el juramento a David; ``no mentiré'', ``testigo fiel en el cielo''};
 \node[draw=gris,line width=0.9pt,fill=gris!12,rounded corners=3pt,text width=3.4cm,align=center,inner sep=5pt,minimum height=1.9cm,right=0.5cm of b] (c) {\textbf{vv. 38--51}\\ Lamento: la corona en el suelo; ``¿dónde están tus antiguas misericordias?''};
 \node[draw=dorado,line width=1.4pt,fill=azul,text=papiro,rounded corners=3pt,text width=2.2cm,align=center,inner sep=5pt,minimum height=1.9cm,right=0.5cm of c] (d) {\textbf{v. 52}\\ ``Bendito sea \SENOR{} para siempre. Amén, y Amén''};
 \draw[-{Latex[length=2.5mm]},dorado,line width=1.2pt] (a)--(b);
 \draw[-{Latex[length=2.5mm]},dorado,line width=1.2pt] (b)--(c);
 \draw[-{Latex[length=2.5mm]},dorado,line width=1.2pt] (c)--(d);
\end{tikzpicture}\par
{\footnotesize\color{gris}\textit{Figura 3.} Las tres partes del Salmo 89 y su doxología. \textit{Nota.} Elaboración propia; numeración española.}
\end{center}

El poeta canta la fidelidad sobre las ruinas de la monarquía davídica. No niega el juramento, sino que se lo recuerda a Dios: es la fe que reclama la fidelidad precisamente cuando no se la ve, la misma que Pink describe: ``Hay épocas en la vida de todas las personas cuando no es fácil, ni siquiera para los cristianos, creer que Dios es fiel'' (Pink, 2020, cap. 10, p. 54). El v. 52 cierra el tercer libro de los Salmos con un ``Amén, y Amén'', palabra de la misma raíz que la fidelidad. La respuesta a la pregunta del v. 49 llegará en Isaías 55:3 y en Pablo (sección~\ref{sec:davidfiel}).

\precision{la división del salmo en himno, oráculo y lamento es la habitual en los comentarios; la ocasión histórica del lamento (la caída de un rey davídico, probablemente en el destierro) se discute. El v. 52 es la doxología que cierra el Libro III del Salterio (Sal 73--89), no una parte original del poema.}

\subsection{“Grande es tu fidelidad”: en el centro del acróstico}
Lamentaciones, escrito sobre las ruinas de Jerusalén, sitúa en su centro la frase más conocida sobre la fidelidad de Dios:

\hebverso{חַסְדֵי יְהוָה כִּי לֹא־תָמְנוּ כִּי לֹא־כָלוּ רַחֲמָיו \\ חֲדָשִׁים לַבְּקָרִים רַבָּה אֱמוּנָתֶךָ}
\begin{verso}{Lamentaciones 3:21–24}
Esto recapacitaré en mi corazón, por lo tanto esperaré. Por la misericordia de \SENOR{} no hemos sido consumidos, porque nunca decayeron sus misericordias. Nuevas son cada mañana; grande es tu fidelidad. Mi porción es \SENOR, dijo mi alma; por tanto, en él esperaré.
\end{verso}

El capítulo 14 estudió estos versos desde la misericordia; aquí interesa su última palabra: \heb{רַבָּה אֱמוּנָתֶךָ} (\emph{rabbá emunatéja}), ``grande es tu fidelidad''. Esa fidelidad se renueva ``cada mañana'', como el maná. El poeta no la deduce de las circunstancias, que son de ruina, sino de la memoria: ``esto recapacitaré en mi corazón'' (\heb{זֹאת אָשִׁיב אֶל־לִבִּי}, ``esto haré volver a mi corazón''). Su respuesta es esperar (\heb{אוֹחִיל}), y lo dice dos veces (vv. 21, 24).

\precision{(textual) el texto griego de la Septuaginta usado en esta serie no contiene Lam 3:22--24, sino que salta del v. 21 al 25 (conviene cotejarlo en la edición de Rahlfs-Hanhart). Los versos están en el texto masorético, que la RVR1960 sigue.}

\subsection{La letra que falta: \emph{nun} en el Salmo 145}
El Salmo 145 es un acróstico en el que cada verso comienza con una letra del alfabeto hebreo. En el texto masorético, sin embargo, falta la letra \emph{nun} (\heb{נ}), que debería ir entre el v. 13 (\emph{mem}) y el v. 14 (\emph{sámej}), y por eso la RVR1960, que sigue ese texto, no tiene el verso. Otros testigos sí lo conservan:

\begin{center}\small\renewcommand{\arraystretch}{1.25}
\begin{tabular}{@{}L{3.6cm}L{10.8cm}@{}}
\toprule
\textbf{Testigo} & \textbf{El verso de \emph{nun}} \\ \midrule
Texto masorético & no lo tiene; el acróstico pasa de \emph{mem} a \emph{sámej} \\
11QPs\textsuperscript{a} (Qumrán) & \heb{נאמן אלוהים בדבריו וחסיד בכול מעשיו}: ``Fiel es Dios en sus palabras, y bondadoso en todas sus obras'' \\
Septuaginta (Sal 144:13a) & \grc{πιστὸς κύριος ἐν τοῖς λόγοις αὐτοῦ καὶ ὅσιος ἐν πᾶσι τοῖς ἔργοις αὐτοῦ}: ``Fiel es el Señor en sus palabras, y santo en todas sus obras'' \\
\bottomrule
\end{tabular}\par
{\footnotesize\color{gris}\textit{Tabla 6.} El verso de \emph{nun} del Salmo 145. \textit{Nota.} 11QPs\textsuperscript{a} según Sanders (1965); LXX según Rahlfs-Hanhart.}
\end{center}

El verso ausente comienza precisamente con \emph{ne'emán}: ``Fiel es Dios en sus palabras''. Y el que le sigue en el texto masorético emplea el verbo del texto ancla de este capítulo: ``Sostiene \SENOR{} a todos los que caen (\heb{הַנֹּפְלִים}, \emph{hanofelim}), y levanta a todos los oprimidos'' (Sal 145:14).

\precision{(textual) los estudiosos están divididos: unos piensan que el verso de \emph{nun} se perdió en la copia del texto masorético; otros, que fue añadido después para completar el acróstico. La RVR1960 sigue el texto masorético. El texto de 11QPs\textsuperscript{a} se cita desde Sanders (1965); conviene cotejarlo con las imágenes del rollo en la Leon Levy Dead Sea Scrolls Digital Library.}

El Talmud ya conocía el salmo sin \emph{nun} y se preguntó por la razón de esa ausencia. Su respuesta se estudia en la sección sobre el judaísmo.

%===================================================================
\section{Teología: el Dios que no puede negarse a sí mismo}

La teología de la fidelidad parte del Nombre que la revela. Moisés no presenta a Israel una cualidad que Dios posea ocasionalmente, sino lo que Dios es: ``Conoce, pues, que \SENOR{} tu Dios es Dios, Dios fiel, que guarda el pacto y la misericordia a los que le aman y guardan sus mandamientos, hasta mil generaciones'' (Dt 7:9). Si la fidelidad pertenece a su ser, Dios no podría dejar de ser fiel sin dejar de ser Dios. Pablo lo afirma con una frase que sorprende justamente porque toma en serio la infidelidad humana: ``Si fuéremos infieles, él permanece fiel; él no puede negarse a sí mismo'' (2~Ti 2:13). La fidelidad de Dios no depende de la respuesta del hombre, porque su raíz no está en el hombre, sino en Dios mismo. Pink lo observa con precisión: ``esta cualidad es esencial para su Ser; sin ella no sería Dios. Para Dios, ser infiel, sería actuar en contra de su naturaleza, lo cual sería imposible'' (Pink, 2020, cap. 10, p. 53); de ahí que ``él nunca olvida, nunca falla, nunca vacila, nunca falta a su Palabra'' (Pink, 2020, cap. 10, pp. 53--54). Es, formulada en negativo, la misma definición de esta serie.

\subsection{La fidelidad en la oscuridad}
Que Dios es fiel se cree con facilidad en la teoría y con dificultad en la prueba. Pink lo advierte: ``una cosa es aceptar la fidelidad de Dios como una verdad divina y otra muy distinta es actuar de acuerdo con ella'' (Pink, 2020, cap. 10, p. 54). Isaías se dirige precisamente al creyente que no ve:

\begin{verso}{Isaías 50:10}
¿Quién hay entre vosotros que teme a \SENOR, y oye la voz de su siervo? El que anda en tinieblas y carece de luz, confíe en el nombre de \SENOR, y apóyese en su Dios.
\end{verso}

``Apóyese'' traduce \heb{וְיִשָּׁעֵן} (\emph{veyisha'én}), ``que se recueste'': quien no ve el camino descansa su peso en Aquel que no falla. El profeta no promete luz inmediata, sino un apoyo firme en medio de la oscuridad. Jesús dirá algo semejante a Pedro: ``Lo que yo hago, tú no lo comprendes ahora; mas lo entenderás después'' (Jn 13:7). La demora, por tanto, no es indiferencia, sino una espera de Dios mismo: ``\SENOR{} esperará para tener piedad de vosotros'' (Is 30:18). La fidelidad de Dios incluye también sus tiempos.

\subsection{La fidelidad en tres obras}
La Escritura muestra esa fidelidad en tres obras que abarcan toda la vida del pueblo de Dios. La primera es preservar. Pablo escribe a los corintios que Cristo ``os confirmará hasta el fin, para que seáis irreprensibles en el día de nuestro Señor Jesucristo. Fiel es Dios, por el cual fuisteis llamados a la comunión con su Hijo'' (1~Co 1:8--9). La perseverancia del creyente no descansa, en palabras de Pink, ``en la fuerza de sus resoluciones'', sino ``en la veracidad de Aquel que no puede mentir'' (Pink, 2020, cap. 10, p. 56).

La segunda obra es disciplinar, y es la que más cuesta reconocer como fidelidad. El salmista, sin embargo, la reconoce: ``Conozco, oh \SENOR, que tus juicios son justos, y que conforme a tu fidelidad me afligiste'' (Sal 119:75). Lejos de contradecir la fidelidad, la aflicción nace de ella, porque Dios no abandona a los suyos a su extravío. Por eso, como observa Pink, Dios ``es tan fiel cuando retiene algo como cuando da algo'' (Pink, 2020, cap. 10, pp. 56--57).

La tercera obra es glorificar: el Dios que comenzó la obra la llevará a término. ``Fiel es el que os llama, el cual también lo hará'' (1~Ts 5:24). Pablo encadena el llamado con la gloria final sin dejar un eslabón suelto: ``a los que llamó, a éstos también justificó; y a los que justificó, a éstos también glorificó'' (Ro 8:30). Preservar, disciplinar y glorificar son, así, tres momentos de una misma fidelidad.

\subsection{Descansar en el fiel Creador}
De esa certeza nace una manera de vivir. La fidelidad de Dios, escribe Pink, ``nos preservará de la preocupación'', ``refrenará nuestros sollozos'' y ``engendrará una confianza cada vez mayor en Dios'' (Pink, 2020, cap. 10, pp. 57--58). Pedro convierte esa certeza en exhortación para los que sufren: ``De modo que los que padecen según la voluntad de Dios, encomienden sus almas al fiel Creador, y hagan el bien'' (1~P 4:19). Encomendar el alma significa entregar en depósito lo más valioso a quien no puede fallar; hacer el bien es la respuesta de quien ya no necesita defenderse a sí mismo.

\precision{Pink (2020, cap. 10) presenta también a Faraón, Coré y Acán como ejemplos de fidelidad en el juicio, y cita un himno de William Cowper (\emph{God moves in a mysterious way}). El capítulo 13 estudió lo que corresponde a la veracidad (Nm 23:19 y las promesas a Abraham y por Isaías); aquí se estudia lo que es propio de la fidelidad.}

%===================================================================
\section{Un encuentro en la historia de Israel: Salomón en la dedicación del templo}

Unos cuatro siglos después de Josué, Salomón concluyó el templo que David había querido edificar. Dios había dicho a David que no lo construiría él, sino su hijo, y le había prometido una casa, es decir, una dinastía (2~S 7:12--16). En la dedicación, con el arca ya en el lugar santísimo y la nube llenando la casa (1~R 8:10--11), Salomón se vuelve hacia el pueblo y lo bendice:

\hebverso{בָּרוּךְ יְהוָה אֱלֹהֵי יִשְׂרָאֵל אֲשֶׁר דִּבֶּר בְּפִיו אֵת דָּוִד אָבִי וּבְיָדוֹ מִלֵּא}
\begin{verso}{1 Reyes 8:15, 20}
Y dijo: Bendito sea \SENOR, Dios de Israel, que habló a David mi padre lo que con su mano ha cumplido, diciendo\ldots{} Y \SENOR{} ha cumplido su palabra que había dicho; porque yo me he levantado en lugar de David mi padre, y me he sentado en el trono de Israel, como \SENOR{} había dicho, y he edificado la casa al nombre de \SENOR{} Dios de Israel.
\end{verso}

\begin{filologia}{1 R 8:15--26}
\heb{דִּבֶּר בְּפִיו} \ldots{} \heb{וּבְיָדוֹ מִלֵּא} & \heb{מלא} & ``habló con su boca\ldots{} y con su mano llenó'' (v. 15; también v. 24). La boca promete; la mano llena la promesa. \\
\heb{וַיָּקֶם יְהוָה אֶת־דְּבָרוֹ} \ldots{} \heb{וָאָקֻם} & \heb{קום} & ``y levantó el \SENOR{} su palabra\ldots{} y yo me levanté'' (v. 20). Salomón mismo es la palabra puesta en pie. \\
\heb{יֵאָמֶן נָא דְּבָרֶיךָ} & \heb{אמן} & ``sea afirmada, te ruego, tu palabra'' (v. 26): nifal de \emph{amán}, la raíz de \emph{ne'emán}. RVR1960: ``cúmplase la palabra''. \\
\heb{לֹא־נָפַל דָּבָר אֶחָד} & \heb{נפל} & ``no cayó una sola palabra'' (v. 56): la frase de Jos 21:45 y 23:14. \\
\end{filologia}

La oración que sigue se construye sobre la fidelidad. Salomón comienza confesando lo que Dios ya hizo y, desde ahí, pide lo que todavía falta:

\begin{verso}{1 Reyes 8:23–26}
\SENOR{} Dios de Israel, no hay Dios como tú, ni arriba en los cielos ni abajo en la tierra, que guardas el pacto y la misericordia a tus siervos, los que andan delante de ti con todo su corazón; que has cumplido a tu siervo David mi padre lo que le prometiste; lo dijiste con tu boca, y con tu mano lo has cumplido, como sucede en este día. Ahora, pues, \SENOR{} Dios de Israel, cumple a tu siervo David mi padre lo que le prometiste\ldots{} Ahora, pues, oh \SENOR{} Dios de Israel, cúmplase la palabra que dijiste a tu siervo David mi padre.
\end{verso}

La fórmula ``que guardas el pacto y la misericordia'' reproduce Dt 7:9 casi palabra por palabra: Salomón ora invocando el Nombre del Dios fiel. Su oración sigue la lógica de toda oración bíblica de confianza: \emph{lo que ya cumpliste es la garantía de lo que te pido que cumplas}. Así había orado también David: ``confirma para siempre la palabra que has hablado sobre tu siervo\ldots{} y haz conforme a lo que has dicho'' (2~S 7:25).

Al terminar, Salomón bendice al pueblo con las palabras de Josué:

\begin{verso}{1 Reyes 8:56–57}
Bendito sea \SENOR, que ha dado paz a su pueblo Israel, conforme a todo lo que él había dicho; ninguna palabra de todas sus promesas que expresó por Moisés su siervo, ha faltado. Esté con nosotros \SENOR{} nuestro Dios, como estuvo con nuestros padres, y no nos desampare ni nos deje.
\end{verso}

``Reposo'' (\heb{מְנוּחָה}) y ``ninguna palabra ha caído'': con el templo se cumple el reposo que Josué había iniciado (Jos 21:44) y que Dt 12:9--10 anunciaba. Dos generaciones de la historia de Israel, la de Josué y la de Salomón, suscriben la misma frase.

El relato, con todo, no termina en la fiesta. En la segunda aparición, Dios advierte a Salomón: ``si obstinadamente os apartareis de mí\ldots{} yo cortaré a Israel de sobre la faz de la tierra\ldots{} y esta casa\ldots{} yo la echaré de delante de mí'' (1~R 9:6--7). Resuena otra vez Jos 23:15, con su palabra buena y su palabra mala. Casi cuatro siglos más tarde, el templo fue quemado (2~R 25:9); y aun entonces la fidelidad siguió hablando: ``Cuando en Babilonia se cumplan los setenta años, yo os visitaré, y despertaré sobre vosotros mi buena palabra'' (Jer 29:10). Por eso el libro de Esdras se abre con el decreto de Ciro, promulgado ``para que se cumpliese la palabra de \SENOR{} por boca de Jeremías'' (Esd 1:1).

\precision{``Despertaré\ldots{} mi buena palabra'' (Jer 29:10) traduce \heb{וַהֲקִמֹתִי עֲלֵיכֶם אֶת־דְּבָרִי הַטּוֹב}, ``levantaré sobre vosotros mi palabra buena'', con el mismo verbo de 1~R 8:20 y la misma ``palabra buena'' de Jos 21:45. La cadena Josué--Salomón--Jeremías es una lectura propia de esta serie, apoyada en ese vocabulario común.}

%===================================================================
\section{La fidelidad en el judaísmo}

\subsection{“Dios, Rey fiel”: el sentido del amén}
El Talmud de Babilonia se pregunta qué significa la palabra con que el pueblo responde a cada bendición:

\hebverso{מַאי אָמֵן אָמַר רַבִּי חֲנִינָא אֵל מֶלֶךְ נֶאֱמָן}

``¿Qué es \emph{amén}? Dijo Rabí Janiná: \emph{El Mélej Ne'emán}'', ``Dios, Rey fiel'' (\emph{Shabat} 119b). Las tres palabras comienzan con las tres letras de \heb{אָמֵן}, álef, mem y nun, de modo que quien dice ``amén'' confiesa que el que bendijo es un Rey que cumple.

\precision{interpretación rabínica (Talmud de Babilonia, edición William Davidson en Sefaria). La lectura de \emph{amén} como sigla es un recurso homilético, no una etimología: \emph{amén} procede de la raíz \emph{amán} (capítulo 13). Con todo, la sigla expresa bien lo que la raíz significa.}

\subsection{“El Dios fiel, que dice y hace”}
Ya en el siglo I se leía en la sinagoga cada sábado, después de la Torá, una sección de los Profetas; fue ``después de la lectura de la ley y de los profetas'' cuando Pablo recibió la invitación a hablar en Antioquía de Pisidia (Hch 13:15). Esa lectura se llama \emph{haftará}, y el tratado \emph{Masejet Soferim} (cap. 13) conserva la bendición que se pronuncia tras ella. En el sidur sefardí reza así:

\hebverso{הָאֵל הַנֶּאֱמָן הָאוֹמֵר וְעוֹשֶׂה הַמְדַבֵּר וּמְקַיֵּם שֶׁכָּל דְּבָרָיו אֱמֶת וָצֶדֶק \\ נֶאֱמָן אַתָּה הוּא יְהוָה אֱלֹהֵינוּ וְנֶאֱמָנִים דְּבָרֶיךָ \\ וְדָבָר אֶחָד מִדְּבָרֶיךָ אָחוֹר לֹא יָשׁוּב רֵיקָם כִּי אֵל מֶלֶךְ נֶאֱמָן וְרַחֲמָן אָתָּה}

\begin{quote}\itshape
\ldots{}el Dios fiel, que dice y hace, que habla y cumple, porque todas sus palabras son verdad y justicia. Fiel eres tú, \SENOR{} nuestro Dios, y fieles son tus palabras, y ni una sola de tus palabras vuelve atrás vacía, porque tú eres Dios, Rey fiel y misericordioso. Bendito eres tú, \SENOR, el Dios fiel en todas sus palabras.\par\hfill{\upshape\small\bfseries\textcolor{azul}{(Bendición después de la \emph{haftará}; \emph{Masejet Soferim} 13)}}
\end{quote}

La bendición ofrece una definición litúrgica del atributo, tejida con textos bíblicos. ``El Dios fiel'' procede de Dt 7:9. ``Que habla y cumple'' (\heb{הַמְדַבֵּר וּמְקַיֵּם}, ``que habla y pone en pie'') emplea el verbo \emph{qum} de 1~R 8:20. ``Ni una sola de tus palabras vuelve atrás vacía'' recoge Is 55:11: ``así será mi palabra que sale de mi boca; no volverá a mí vacía'' (\heb{לֹא־יָשׁוּב אֵלַי רֵיקָם}). El cierre, por último, ``Dios, Rey fiel'', es la sigla del amén de \emph{Shabat} 119b.

\precision{texto según el \emph{Sidur Sefard} (servicio de la mañana del sábado) y \emph{Masejet Soferim} 13 en Sefaria; las variantes entre sidurim son menores. Sobre la historia de la lectura de los Profetas y sus bendiciones, véase Elbogen (1993). No es casual que la bendición se pronuncie justo después de los Profetas: en ellos se anuncian las promesas que Dios cumplirá.}

\subsection{La \emph{Amidá}: fiel a los que duermen en el polvo}
La segunda bendición de la \emph{Amidá}, la oración que el judío recita tres veces al día, alaba el poder de Dios (capítulo 9), y en ella aparecen dos frases sobre la fidelidad:

\hebverso{וּמְקַיֵּם אֱמוּנָתוֹ לִישֵׁנֵי עָפָר \\ וְנֶאֱמָן אַתָּה לְהַחֲיוֹת מֵתִים}

``\ldots{}y cumple su fidelidad a los que duermen en el polvo'' y ``fiel eres tú para dar vida a los muertos''. La primera toma la expresión de Dn 12:2: ``muchos de los que duermen en el polvo de la tierra serán despertados''. La fidelidad de Dios, entonces, no termina en la tumba: la promesa hecha a los padres se cumplirá también en los muertos. En la misma bendición, Dios es \heb{סוֹמֵךְ נוֹפְלִים}, ``el que sostiene a los que caen'', con palabras de Sal 145:14.

\precision{texto según el \emph{Sidur Ashkenaz} en Sefaria. La \emph{Amidá} se fijó en su forma básica después de la destrucción del segundo templo, y la redacción exacta de algunas frases varió durante siglos (Elbogen, 1993). La conexión entre esta fidelidad ``a los que duermen'' y la resurrección es explícita en la bendición misma.}

\subsection{\emph{Modé aní}: la fidelidad cada mañana}
La primera oración del día, que se dice al despertar, todavía en la cama, concluye con las palabras de Lamentaciones:

\hebverso{מוֹדֶה אֲנִי לְפָנֶיךָ מֶלֶךְ חַי וְקַיָּם \\ שֶׁהֶחֱזַרְתָּ בִּי נִשְׁמָתִי בְּחֶמְלָה רַבָּה אֱמוּנָתֶךָ}

``Te doy gracias, Rey vivo y permanente, porque me devolviste el alma con compasión: grande es tu fidelidad''. La oración lee Lam 3:23 al pie de la letra, pues las misericordias ``nuevas cada mañana'' son el aliento que vuelve al despertar.

\precision{texto según el \emph{Sidur Ashkenaz} en Sefaria. \emph{Modé aní} es una oración tardía, que aparece por primera vez en el \emph{Séder haYom} de Moshé ben Majir (s.~XVI). La idea de que el alma se entrega a Dios al dormir y se recibe de vuelta al despertar es, en cambio, más antigua en la tradición judía.}

\subsection{La letra que falta y el que sostiene a los que caen}
El Talmud conocía el Salmo 145 sin el verso de \emph{nun} y propuso una explicación:

\begin{quote}\itshape
Dijo Rabí Yojanán: ¿Por qué no se dice \emph{nun} en el \emph{Ashré} (Sal 145)? Porque en ella está la caída de [los enemigos de] Israel, como está escrito: ``Cayó la virgen de Israel, y no podrá levantarse ya más'' (Am 5:2)\ldots{} Dijo Rav Najmán bar Yitsjak: Aun así, David volvió a sostenerlos por el Espíritu Santo, como está dicho: ``Sostiene el \SENOR{} a todos los que caen'' (Sal 145:14).\par\hfill{\upshape\small\bfseries\textcolor{azul}{(\emph{Berajot} 4b)}}
\end{quote}

\hebverso{סוֹמֵךְ יְהוָה לְכָל־הַנֹּפְלִים}

Am 5:2 comienza con \heb{נָפְלָה} (\emph{nafelá}), ``cayó''. Según el Talmud, David no quiso iniciar un verso con la letra de la caída; pero en el verso siguiente, que empieza con \emph{sámej}, el \SENOR{} sostiene a ``los que caen''. Se trata del mismo verbo de Jos 21:45: la palabra de Dios no cae, y Dios levanta a los que caen.

\precision{interpretación rabínica; ``los enemigos de Israel'' es un eufemismo del Talmud para no decir ``la caída de Israel''. La lectura que une \emph{Berajot} 4b con el verbo de Jos 21:45 es propia de esta serie. El verso de \emph{nun} que conservan Qumrán y la LXX (``Fiel es Dios en sus palabras'') no forma parte del texto que comentaba el Talmud.}

\subsection{“Aunque demore su recompensa”}
El Sifré, un midrash antiguo sobre Deuteronomio, comenta el título \emph{El emuná} de Dt 32:4 con una frase atrevida:

\hebverso{אֵל אֱמוּנָה שֶׁהֶאֱמִין בָּעוֹלָם וּבְרָאוֹ}

``Dios de fidelidad: que tuvo confianza en el mundo y lo creó'' (\emph{Sifré Devarim} 307). El midrash juega con la raíz, pues Dios es fiel y, a la vez, ``confió'' al crear; y añade que Dios retribuye a cada uno con exactitud, en este mundo o en el venidero. Rashi resume el pasaje con una frase que toca el punto más difícil de la fidelidad, la demora:

\begin{quote}\itshape
Para pagar a los justos su justicia en el mundo venidero; y aunque demore su recompensa, al final hace firmes sus palabras.\par\hfill{\upshape\small\bfseries\textcolor{azul}{(Rashi a Dt 32:4)}}
\end{quote}

``Al final hace firmes sus palabras'' traduce \heb{סוֹפוֹ לְאַמֵּן אֶת דְּבָרָיו}, y el verbo es, una vez más, \emph{amán}. La fidelidad abarca también el tiempo de espera.

\precision{interpretaciones rabínicas, según Sefaria (\emph{Sifré Devarim} 307; Rashi a Dt 32:4).}

\subsection{Qumrán: la espera y la fidelidad}
El \emph{Pésher de Habacuc} (1QpHab), uno de los primeros rollos hallados en Qumrán, comenta Hab 2:3--4 desde la experiencia de una comunidad que esperaba el fin y constataba que no llegaba. A propósito de ``aunque la visión tardará aún por un tiempo\ldots{} aunque tardare, espéralo'' (Hab 2:3), explica que ``el tiempo final se prolongará y excederá todo lo que dijeron los profetas, porque los misterios de Dios son asombrosos'', y que el texto alude a ``los hombres de verdad que cumplen la Ley, cuyas manos no se aflojarán en el servicio de la verdad cuando el tiempo final se prolongue'' (1QpHab VII 5--14). En cuanto a ``el justo por su fe vivirá'' (Hab 2:4), lo refiere a quienes cumplen la Ley en la casa de Judá, a los que Dios librará del juicio ``por su sufrimiento y por su fidelidad al Maestro de Justicia'' (1QpHab VIII 1--3).

\precision{traducción propia desde la versión de Vermes (2004); cf. García Martínez (1992). El término traducido ``fidelidad'' es la \emph{emuná} de Hab 2:4. Que Qumrán y Pablo comenten el mismo versículo, con lecturas distintas, muestra que Hab 2:4 era un texto clave del judaísmo del siglo I.}

%===================================================================
\section{La fidelidad en Cristo, leída con Pablo}

\subsection{“No que la palabra de Dios haya fallado”}
En Romanos 9--11 Pablo enfrenta la pregunta más dura que un judío de su tiempo podía plantear sobre la fidelidad de Dios: si el Mesías vino y la mayor parte de Israel no lo recibió, ¿falló la promesa? Tras enumerar los privilegios de Israel, ``la adopción, la gloria, el pacto, la promulgación de la ley, el culto y las promesas'' (Ro 9:4), responde:

\begin{verso}{Romanos 9:6}
No que la palabra de Dios haya fallado; porque no todos los que descienden de Israel son israelitas.
\end{verso}

\begin{filologia}{Ro 9:6}
\grc{οὐχ οἷον δὲ ὅτι ἐκπέπτωκεν} & \grc{ἐκπίπτω} & ``no es que haya caído'': perfecto de \grc{ἐκπίπτω}, ``caer fuera, caer al suelo''. La RVR1960 traduce ``haya fallado''. \\
\grc{ὁ λόγος τοῦ θεοῦ} & \grc{λόγος} & ``la palabra de Dios'': las promesas a los padres (Ro 9:4, 8--9). \\
\grc{οὐ διέπεσεν} (Jos 21:45 LXX) & \grc{πίπτω} & la Septuaginta de Josué usa la misma imagen: la palabra ``no cayó''. \\
\end{filologia}

Pablo formula la pregunta del Salmo 89:49 y responde con la certeza de Josué: la palabra no cayó. Los capítulos 9--11 desarrollan esa respuesta hasta llegar a ``¿Ha desechado Dios a su pueblo? En ninguna manera'' (Ro 11:1) y a ``irrevocables son los dones y el llamamiento de Dios'' (Ro 11:29; capítulo 8). Ya en Ro 3:3 había planteado la misma cuestión: ``¿Su incredulidad habrá hecho nula la fidelidad de Dios?'' (\grc{τὴν πίστιν τοῦ θεοῦ}; capítulo 13).

\precision{la relación verbal entre Ro 9:6 (\grc{ἐκπέπτωκεν}) y Jos 21:45 LXX (\grc{οὐ διέπεσεν}) es una lectura propia de esta serie: los verbos pertenecen a la misma familia y la imagen es la misma, pero Pablo no cita Josué. Que la pregunta de Ro 9--11 sea la fidelidad de Dios a Israel es la lectura común de los comentarios (Moo, 1996; Wright, 2013).}

\subsection{“Fiel es Dios”: una fórmula de Pablo}
Pablo convierte el título de Dt 7:9 en una fórmula que repite a lo largo de sus cartas, y siempre en el mismo punto del argumento: cuando la iglesia mira al futuro y teme no llegar.

\begin{center}\small\renewcommand{\arraystretch}{1.22}
\begin{tabular}{@{}L{2.2cm}L{5.4cm}L{6.8cm}@{}}
\toprule
\textbf{Texto} & \textbf{Griego} & \textbf{Lo que Dios hará por ser fiel} \\ \midrule
1~Co 1:9 & \grc{πιστὸς ὁ θεός} & confirmar a los creyentes ``hasta el fin'' (v. 8) en la comunión con su Hijo \\
1~Co 10:13 & \grc{πιστὸς δὲ ὁ θεός} & no permitir la prueba más allá de lo que se puede resistir; dar la salida \\
2~Co 1:18 & \grc{πιστὸς δὲ ὁ θεός} & que su palabra no sea ``Sí y No'' (capítulo 13) \\
1~Ts 5:24 & \grc{πιστὸς ὁ καλῶν ὑμᾶς} & santificar por completo y guardar hasta la venida de Cristo \\
2~Ts 3:3 & \grc{πιστὸς δέ ἐστιν ὁ κύριος} & afirmar y guardar del mal \\
2~Ti 2:13 & \grc{ἐκεῖνος πιστὸς μένει} & permanecer fiel aunque seamos infieles \\
\bottomrule
\end{tabular}\par
{\footnotesize\color{gris}\textit{Tabla 7.} La fórmula ``Fiel es Dios'' en las cartas de Pablo. \textit{Nota.} Texto griego NA28.}
\end{center}

\precision{la autoría paulina de 2 Tesalonicenses y 2 Timoteo se discute. La fórmula, sin embargo, está en cartas indiscutidas (1 Corintios, 2 Corintios, 1 Tesalonicenses), y la idea de 2~Ti 2:13 se encuentra ya en Ro 3:3--4.}

\subsection{Llamados a la comunión del Hijo}
\begin{verso}{1 Corintios 1:8–9}
el cual también os confirmará hasta el fin, para que seáis irreprensibles en el día de nuestro Señor Jesucristo. Fiel es Dios, por el cual fuisteis llamados a la comunión con su Hijo Jesucristo nuestro Señor.
\end{verso}

\begin{filologia}{1 Co 1:8--9}
\grc{βεβαιώσει ὑμᾶς ἕως τέλους} & \grc{βεβαιόω} & ``os hará firmes hasta el fin'': el verbo griego de la firmeza, cercano a \emph{amán}. \\
\grc{πιστὸς ὁ θεός} & \grc{πιστός} & ``fiel es Dios'': sin verbo, como una confesión. LXX Dt 7:9: \grc{θεὸς πιστός}. \\
\grc{εἰς κοινωνίαν τοῦ υἱοῦ αὐτοῦ} & \grc{κοινωνία} & ``a la comunión de su Hijo'': participar en él, ser de él. \\
\end{filologia}

Pablo se dirige aquí a la iglesia más problemática de sus cartas, marcada por divisiones, inmoralidad, pleitos y desorden en la Cena. Antes de corregir, sin embargo, establece el fundamento: los corintios llegarán al día de Cristo no por su propia firmeza, sino porque Dios es fiel. Pink subraya el mismo punto: la confianza del apóstol ``no se basó en la fuerza de sus resoluciones o capacidad de perseverar, sino en la veracidad de Aquel que no puede mentir'' (Pink, 2020, cap. 10, p. 56).

\subsection{Fiel en la prueba: la roca que los seguía}
En 1 Corintios 10 Pablo relee la historia del desierto. Israel pasó el mar, comió el maná y ``bebían de la roca espiritual que los seguía, y la roca era Cristo'' (10:4); pero cayó en idolatría, fornicación y murmuración, y ``cayeron en un día veintitrés mil'' (10:8). De ahí la advertencia: ``el que piensa estar firme, mire que no caiga'' (10:12). Y a continuación, el consuelo:

\begin{verso}{1 Corintios 10:13}
No os ha sobrevenido ninguna tentación que no sea humana; pero fiel es Dios, que no os dejará ser tentados más de lo que podéis resistir, sino que dará también juntamente con la tentación la salida, para que podáis soportar.
\end{verso}

El contraste es el que recorre todo este capítulo: el pueblo cae, Dios no. El Cántico de Moisés había llamado a Dios ``la Roca'' y \emph{El emuná}, título que la Septuaginta traduce \grc{θεὸς πιστός} (Dt 32:4); Pablo reúne ambas ideas en un mismo pasaje, donde Cristo es la roca que acompaña y Dios es fiel en la prueba.

\precision{es probable que Pablo tenga en mente Dt 32 en 1~Co 10, pues cita Dt 32:17 en 10:20 y alude a Dt 32:21 en 10:22 (Fee, 1987; Hays, 1989). El vínculo entre ``la roca'' de 10:4 y \grc{θεὸς πιστός} de 10:13 a través de Dt 32:4 es una lectura propia.}

\subsection{“El cual también lo hará”}
\begin{verso}{1 Tesalonicenses 5:23–24}
Y el mismo Dios de paz os santifique por completo; y todo vuestro ser, espíritu, alma y cuerpo, sea guardado irreprensible para la venida de nuestro Señor Jesucristo. Fiel es el que os llama, el cual también lo hará.
\end{verso}

En la frase \grc{πιστὸς ὁ καλῶν ὑμᾶς, ὃς καὶ ποιήσει} resuena, en boca de Pablo, la definición de esta serie: el que llama, cumple. La misma certeza reaparece en una carta indiscutida escrita desde la cárcel:

\begin{verso}{Filipenses 1:6}
estando persuadido de esto, que el que comenzó en vosotros la buena obra, la perfeccionará hasta el día de Jesucristo.
\end{verso}

``El que comenzó\ldots{} la perfeccionará'' (\grc{ὁ ἐναρξάμενος\ldots{} ἐπιτελέσει}): Dios no deja obras a medio hacer. Así creyó Abraham, ``plenamente convencido de que era también poderoso para hacer todo lo que había prometido'' (Ro 4:21; capítulo 9).

\subsection{“El justo por la fe vivirá”}
Habacuc se queja de que Dios tarda, y Dios le responde con una visión que ``no mentirá'' y que ``sin duda vendrá'' (Hab 2:3). En ese contexto de espera se encuentra la frase que Pablo sitúa al comienzo de Romanos:

\hebverso{וְצַדִּיק בֶּאֱמוּנָתוֹ יִחְיֶה}
\begin{verso}{Habacuc 2:4}
He aquí que aquel cuya alma no es recta, se enorgullece; mas el justo por su fe vivirá.
\end{verso}
\begin{verso}{Romanos 1:17}
Porque en el evangelio la justicia de Dios se revela por fe y para fe, como está escrito: Mas el justo por la fe vivirá.
\end{verso}

\begin{center}\small\renewcommand{\arraystretch}{1.25}
\begin{tabular}{@{}L{2.6cm}L{6.2cm}L{5.6cm}@{}}
\toprule
\textbf{Testigo} & \textbf{Texto} & \textbf{De quién es la fidelidad} \\ \midrule
Hab 2:4 TM & \heb{וְצַדִּיק בֶּאֱמוּנָתוֹ יִחְיֶה} & ``por \emph{su} \emph{emuná}'': la del justo, que se mantiene firme en la espera \\
Hab 2:4 LXX & \grc{ὁ δὲ δίκαιος ἐκ πίστεώς μου ζήσεται} & ``por \emph{mi} fidelidad'': la de Dios \\
1QpHab VIII 1--3 & ``por su fidelidad al Maestro de Justicia'' & la de la comunidad hacia su maestro \\
Ro 1:17; Gá 3:11 & \grc{ὁ δὲ δίκαιος ἐκ πίστεως ζήσεται} & sin pronombre: la fe del que cree, fundada en la fidelidad de Dios \\
He 10:38 & \grc{ὁ δὲ δίκαιός μου ἐκ πίστεως ζήσεται} & ``mi justo'': el pronombre pasa al justo \\
\bottomrule
\end{tabular}\par
{\footnotesize\color{gris}\textit{Tabla 8.} Hab 2:4 en el texto masorético, la LXX, Qumrán y el Nuevo Testamento. \textit{Nota.} Algunos manuscritos de la LXX ponen el pronombre como en Hebreos.}
\end{center}

Los testigos muestran que la frase, como la raíz, tiene dos caras. En hebreo, \emph{emuná} es la firmeza del justo que espera; en la Septuaginta, la fidelidad de Dios que no falla. Pablo deja la palabra sin pronombre y la abre a ambos sentidos: el justo vive por la fe, y esa fe se apoya en un Dios fiel. La fe es, así, la respuesta humana a la fidelidad divina, del mismo modo que \emph{he'emín} (creer) significa apoyarse en lo que es \emph{ne'emán} (firme).

\precision{se discute si el ``por fe'' de Ro 1:17 apunta a la fe del creyente, a la fidelidad de Dios o a la de Cristo. Moo (1996) lo lee como la fe del creyente; Dunn (1988) entiende ``de fe y para fe'' como ``de la fidelidad de Dios a la fe del hombre''; Hays (2002) lo relaciona con la fidelidad de Cristo. La lectura que aquí se ofrece, que une la fe humana a la fidelidad de Dios, está próxima a la de Dunn y no es unánime.}

\subsection{La fidelidad de Cristo}
Pablo emplea varias veces la expresión griega \grc{πίστις Ἰησοῦ Χριστοῦ}, literalmente ``fe de Jesucristo'' (Ro 3:22, 26; Gá 2:16, 20; 3:22; Flp 3:9). La RVR1960 la traduce ``fe en Jesucristo'' en Ro 3:22 y Gá 3:22, pero conserva la forma literal en Gá 2:16 (``por la fe de Jesucristo''), en Ro 3:26 (``el que es de la fe de Jesús'') y en Gá 2:20 (``lo vivo en la fe del Hijo de Dios, el cual me amó y se entregó a sí mismo por mí''). El debate gira en torno a si Pablo habla de la fe que el creyente pone en Cristo o de la fidelidad que Cristo mismo mantuvo hacia Dios.

\begin{center}\small\renewcommand{\arraystretch}{1.25}
\begin{tabular}{@{}L{3.6cm}L{5.4cm}L{5.4cm}@{}}
\toprule
 & \textbf{``Fe en Cristo''} (genitivo objetivo) & \textbf{``Fidelidad de Cristo''} (genitivo subjetivo) \\ \midrule
Lectura & la confianza del creyente puesta en Cristo & la obediencia fiel de Cristo hasta la cruz \\
Argumentos & Pablo usa el verbo ``creer en Cristo'' en el mismo versículo: ``nosotros también hemos creído en Jesucristo'' (Gá 2:16) & Ro 3:3 usa la misma construcción para ``la fidelidad de Dios''; Flp 2:8: ``obediente hasta la muerte'' \\
Autores & Dunn (1998) & Hays (2002); Wright (2013) \\
\bottomrule
\end{tabular}\par
{\footnotesize\color{gris}\textit{Tabla 9.} El debate sobre \grc{πίστις Χριστοῦ}. \textit{Nota.} Elaboración propia.}
\end{center}

Ambas lecturas afirman verdades paulinas: el creyente confía en Cristo, y Cristo fue fiel hasta la muerte. Para este capítulo basta advertir que, en Pablo, la fidelidad de Dios no queda reducida a un atributo abstracto, sino que toma la forma de un hombre ``obediente hasta la muerte, y muerte de cruz'' (Flp 2:8). Donde Israel cayó, Cristo permaneció fiel.

\precision{el debate sigue abierto en los estudios paulinos, y aquí se presentan las dos lecturas sin zanjarlo. La tabla resume posiciones generales; cada autor matiza la suya.}

\subsection{Las misericordias fieles de David}\label{sec:davidfiel}
El Salmo 89 terminaba con una pregunta: ``Señor, ¿dónde están tus antiguas misericordias, que juraste a David por tu verdad?'' (Sal 89:49). Isaías respondió con la promesa de un pacto eterno:

\hebverso{וְאֶכְרְתָה לָכֶם בְּרִית עוֹלָם חַסְדֵי דָוִד הַנֶּאֱמָנִים}
\begin{verso}{Isaías 55:3}
Inclinad vuestro oído, y venid a mí; oíd, y vivirá vuestra alma; y haré con vosotros pacto eterno, las misericordias firmes a David.
\end{verso}

En hebreo, ``las misericordias firmes a David'' son \emph{jasdé David hane'emanim}, las \emph{jésed} de David, las \emph{fieles}: las dos palabras del Salmo 89 reunidas en una sola frase. Y Pablo, según Hechos, predicó en la sinagoga de Antioquía de Pisidia que esa promesa se había cumplido en la resurrección de Jesús:

\begin{verso}{Hechos 13:32–34}
Y nosotros también os anunciamos el evangelio de aquella promesa hecha a nuestros padres, la cual Dios ha cumplido a los hijos de ellos, a nosotros, resucitando a Jesús\ldots{} Y en cuanto a que le levantó de los muertos para nunca más volver a corrupción, lo dijo así: Os daré las misericordias fieles de David.
\end{verso}

\begin{center}
\begin{tikzpicture}[font=\small,node distance=0.7cm]
 \node[draw=gris,line width=0.8pt,fill=gris!12,rounded corners=3pt,text width=3.2cm,align=center,inner sep=6pt] (a) {\textbf{Sal 89:49}\\ ``¿dónde están tus antiguas misericordias?''\\ \heb{חֲסָדֶיךָ} \ldots{} \heb{בֶּאֱמוּנָתֶךָ}};
 \node[draw=azul,line width=0.8pt,fill=papiro,rounded corners=3pt,text width=3.2cm,align=center,inner sep=6pt,right=of a] (b) {\textbf{Is 55:3}\\ ``las misericordias firmes a David''\\ \heb{חַסְדֵי דָוִד הַנֶּאֱמָנִים}};
 \node[draw=dorado,line width=1.6pt,fill=azul,text=papiro,rounded corners=3pt,text width=3.4cm,align=center,inner sep=6pt,right=of b] (c) {\textbf{Hch 13:34}\\ ``os daré las misericordias fieles de David''\\ \grc{τὰ ὅσια Δαυὶδ τὰ πιστά}};
 \draw[-{Latex[length=3mm]},dorado,line width=1.4pt] (a)--(b);
 \draw[-{Latex[length=3mm]},dorado,line width=1.4pt] (b)--(c);
\end{tikzpicture}\par
{\footnotesize\color{gris}\textit{Figura 4.} De la pregunta del Salmo 89 a la resurrección. \textit{Nota.} Elaboración propia.}
\end{center}

La respuesta de Dios a la corona caída de David no consistió en restaurar un trono en Jerusalén, sino en levantar de los muertos al Hijo de David. El verbo con que Hechos designa la resurrección, \grc{ἀνέστησεν} (``levantó'', 13:34), es el mismo con que la Septuaginta traduce 1~R 8:20: \grc{καὶ ἀνέστησεν κύριος τὸ ῥῆμα αὐτοῦ}, ``y el Señor levantó su palabra''. La fidelidad de Dios a David atravesó la tumba y salió de ella.

\precision{el discurso de Hch 13:16--41 es un resumen de Lucas, no una carta de Pablo; pero la idea es paulina, pues Cristo es ``del linaje de David según la carne'' y fue ``declarado Hijo de Dios con poder\ldots{} por la resurrección de entre los muertos'' (Ro 1:3--4). La coincidencia del verbo con 1~R 8:20 LXX está en el texto griego; que Lucas o Pablo la tuvieran en mente es una lectura propia.}

\subsection{El testigo fiel}
El Apocalipsis llama a Jesucristo ``el testigo fiel, el primogénito de los muertos, y el soberano de los reyes de la tierra'' (Ap 1:5), y los tres títulos proceden del Salmo 89: Dios pondrá a David ``por primogénito, el más excelso de los reyes de la tierra'' (89:27), y su trono será ``como un testigo fiel en el cielo'' (89:37; LXX \grc{ὁ μάρτυς ἐν οὐρανῷ πιστός}). Lo que el salmo prometió al rey de Israel, Juan lo contempla cumplido en el Resucitado. Por eso el último libro de la Biblia lo llama ``el Amén, el testigo fiel y verdadero'' (Ap 3:14) y ``Fiel y Verdadero'' (Ap 19:11).

\precision{Juan no es Pablo; se incluye aquí porque cierra el hilo del Salmo 89. Los comentarios señalan el uso del Salmo 89 en Ap 1:5 (Beale, 1999; Bauckham, 1993). Hebreos añade que Cristo ``es fiel al que le constituyó, como también lo fue Moisés en toda la casa de Dios'' (He 3:2), con la frase de Nm 12:7.}

Pablo resume la cristología del atributo en la carta que el capítulo 13 ya estudió: ``todas las promesas de Dios son en él Sí, y en él Amén'' (2~Co 1:20). La fidelidad de Dios no se reduce a que cumpla sus promesas; en Cristo, en quien ``habita corporalmente toda la plenitud de la Deidad'' (Col 2:9), la promesa y el cumplimiento son una misma Persona.

%===================================================================
\section{Aplicación: un atributo comunicable}

\subsection{Hallado fiel}
Dios es fiel en sentido absoluto; el ser humano lo es como reflejo. Por eso la Escritura aplica a algunos hombres la misma palabra que reserva para Dios. De Abraham dicen los levitas: ``hallaste fiel su corazón delante de ti\ldots{} y cumpliste tu palabra, porque eres justo'' (Neh 9:8). El verso pone frente a frente las dos fidelidades: Abraham fue \emph{ne'emán}, y Dios ``levantó'' su palabra (\heb{וַתָּקֶם אֶת־דְּבָרֶיךָ}). De Moisés, Dios mismo afirma que es ``fiel en toda mi casa'' (Nm 12:7); de Samuel se dice que fue ``fiel profeta'' (1~S 3:20), y antes, que Dios ``no dejó caer a tierra ninguna de sus palabras'' (1~S 3:19).

Pablo exige lo mismo a los ministros: ``se requiere de los administradores, que cada uno sea hallado fiel'' (1~Co 4:2). Además, sitúa la fidelidad entre los frutos del Espíritu: \grc{πίστις} en Gá 5:22, que la RVR1960 traduce ``fe'' pero que, en ese contexto de virtudes de trato (paciencia, benignidad, bondad, mansedumbre), significa ``fidelidad'' o ``lealtad''. Jesús lo resume en una sentencia: ``El que es fiel en lo muy poco, también en lo más es fiel'' (Lc 16:10).

\precision{que \grc{πίστις} en Gá 5:22 signifique ``fidelidad'' y no ``fe'' es la lectura de muchos comentarios y de versiones como LBLA (``fidelidad''); la RVR1960 traduce ``fe''.}

\subsection{La infidelidad humana y la fidelidad de Dios}
Precisamente porque es reflejo, la fidelidad humana puede fallar. La Biblia describe la infidelidad de Israel con el lenguaje del adulterio: ``como la esposa infiel abandona a su compañero, así prevaricasteis contra mí'' (Jer 3:20). La respuesta de Dios, sin embargo, no es el abandono: ``Y te desposaré conmigo en fidelidad (\heb{בֶּאֱמוּנָה}), y conocerás a \SENOR'' (Os 2:20 [2:22]). Dios vuelve a desposar a la infiel, y la dote que aporta es su propia fidelidad. Sobre ese fondo escribe Pablo:

\begin{verso}{2 Timoteo 2:13}
Si fuéremos infieles, él permanece fiel; El no puede negarse a sí mismo.
\end{verso}

\subsection{Cinco respuestas}
\begin{enumerate}
\item \textbf{Recordar lo cumplido.} Josué pidió a Israel que ``reconociera'' lo que ya había visto (Jos 23:14), y el poeta de Lamentaciones pudo esperar porque hizo ``volver a su corazón'' la fidelidad de Dios (Lam 3:21--23). La memoria de las promesas cumplidas sostiene la espera de las que aún faltan.
\item \textbf{Orar con sus palabras.} David y Salomón oraron recordando a Dios lo que él mismo había dicho: ``haz conforme a lo que has dicho'' (2~S 7:25); ``cúmplase la palabra'' (1~R 8:26). La oración más firme es la que se apoya en una promesa.
\item \textbf{Esperar en la demora.} ``Aunque tardare, espéralo'' (Hab 2:3); ``aunque demore su recompensa, al final hace firmes sus palabras'' (Rashi a Dt 32:4). De ahí la exhortación de Hebreos: ``Mantengamos firme, sin fluctuar, la profesión de nuestra esperanza, porque fiel es el que prometió'' (He 10:23).
\item \textbf{Apoyarse en él en la oscuridad.} ``El que anda en tinieblas y carece de luz, confíe en el nombre de \SENOR, y apóyese en su Dios'' (Is 50:10). Y en la prueba, ``fiel es Dios, que no os dejará ser tentados más de lo que podéis resistir'' (1~Co 10:13).
\item \textbf{Ser hallado fiel.} La fidelidad se pide en lo poco y en lo mucho (Lc 16:10), en la palabra dada y en el pacto matrimonial, en el servicio y en la confesión: ``Si confesamos nuestros pecados, él es fiel y justo para perdonar nuestros pecados'' (1~Jn 1:9). ``Sé fiel hasta la muerte, y yo te daré la corona de la vida'' (Ap 2:10).
\end{enumerate}

\begin{verso}{Hebreos 10:23}
Mantengamos firme, sin fluctuar, la profesión de nuestra esperanza, porque fiel es el que prometió.
\end{verso}

El atributo deja planteada la pregunta de Jeremías: ¿tratamos a Dios como un arroyo que puede secarse, o como el manantial cuyas aguas no mienten? Y deja también una segunda pregunta, dirigida al creyente mismo: ¿somos, para quienes dependen de nosotros, arroyo o manantial?

%===================================================================
\section{Síntesis}

\begin{tcolorbox}[colback=papiro,colframe=azul,boxrule=1pt,arc=3pt]
Dios cumple todo lo que promete; nunca falla. Josué lo expresó con una imagen: ``no cayó una palabra'' de todas las buenas promesas (Jos 21:45). Su Nombre es \emph{ha'El hane'emán}, ``el Dios fiel'', que guarda el pacto y la \emph{jésed} por mil generaciones (Dt 7:9). Frente al arroyo engañoso, él es el manantial cuyas aguas no mienten (Is 58:11; Jer 15:18). Cumple tanto la palabra buena como la advertencia, espera y enseña a esperar, y permanece fiel aun ante la infidelidad de su pueblo. Israel lo cantó en el Salmo 89 y lo confiesa cada vez que dice ``amén'' (\emph{El Mélej Ne'emán}); Pablo lo resumió en una fórmula, ``Fiel es Dios'' (1~Co 1:9), y lo vio cumplido en Cristo, el Hijo de David resucitado, en quien todas las promesas son ``Sí'' y ``Amén'' (Hch 13:34; 2~Co 1:20).
\end{tcolorbox}

\begin{center}\small\renewcommand{\arraystretch}{1.3}
\begin{tabular}{@{}L{3.0cm}L{11.4cm}@{}}
\toprule
\textbf{Mirada} & \textbf{Aporte al atributo} \\ \midrule
Hebrea & \emph{ne'emán}, \emph{emuná}, \emph{jésed ve'emuná}; los verbos de cumplir (no caer, llenar, levantar); Jos 21:45; Dt 7:9; el arroyo y el manantial; el Salmo 89; Lam 3:23; Salomón en 1~R 8. \\
Judía & \emph{amén} como \emph{El Mélej Ne'emán} (\emph{Shabat} 119b); ``el Dios fiel, que dice y hace'' (bendición de la \emph{haftará}); fiel a los que duermen en el polvo (\emph{Amidá}); \emph{Modé aní}; \emph{Berajot} 4b; Sifré y Rashi sobre la demora; 1QpHab. \\
Teológica & El Dios que no puede negarse a sí mismo (Dt 7:9; 2~Ti 2:13); fiel en la oscuridad (Is 50:10); demostrada al preservar, disciplinar y glorificar (1~Co 1:9; Sal 119:75; 1~Ts 5:24); vivida como descanso en el fiel Creador (1~P 4:19). \\
Paulina & ``Fiel es Dios'' (1~Co 1:9; 10:13; 1~Ts 5:24); la palabra que no cayó (Ro 9:6); el justo por la fe (Ro 1:17); la fidelidad de Cristo (Ro 3:22; Flp 2:8); las misericordias fieles de David (Hch 13:34); el Sí y el Amén (2~Co 1:20). \\
\bottomrule
\end{tabular}\end{center}

\begin{verso}{Jeremías 9:24 [9:23]}
Mas alábese en esto el que se hubiere de alabar: en entenderme y conocerme, que yo soy \SENOR{}, que hago misericordia, juicio y justicia en la tierra; porque estas cosas quiero, dice \SENOR{}.
\end{verso}

``Conoce, pues, que \SENOR{} tu Dios es Dios, Dios fiel'' (Dt 7:9). Conocer a Dios, como pide Jeremías, significa saber por experiencia lo que Josué supo al final de su vida: que ninguna de sus palabras cayó al suelo.

%===================================================================
\section*{Glosario del capítulo}
\begin{description}\setlength{\itemsep}{2pt}
\item[\emph{Ajzav}] (\heb{אַכְזָב}). ``Engañoso'': el arroyo que se seca en verano (Jer 15:18; Mi 1:14).
\item[\emph{Amén}] (\heb{אָמֵן}). ``Así es'', ``firme''; en el Talmud, sigla de \emph{El Mélej Ne'emán} (\emph{Shabat} 119b).
\item[\emph{Bagad}] (\heb{בָּגַד}). Traicionar, ser infiel (Jer 3:20; Job 6:15).
\item[\emph{Ekpíptō}] (\grc{ἐκπίπτω}). Caer, fallar: ``no que la palabra de Dios haya fallado'' (Ro 9:6).
\item[\emph{El emuná}] (\heb{אֵל אֱמוּנָה}). ``Dios de fidelidad'' (Dt 32:4); LXX \grc{θεὸς πιστός}.
\item[\emph{El Mélej Ne'emán}] (\heb{אֵל מֶלֶךְ נֶאֱמָן}). ``Dios, Rey fiel'': fórmula litúrgica judía.
\item[\emph{Emuná}] (\heb{אֱמוּנָה}). Fidelidad, firmeza; también fe (Lam 3:23; Hab 2:4).
\item[\emph{Ha'El hane'emán}] (\heb{הָאֵל הַנֶּאֱמָן}). ``El Dios fiel'': el Nombre del atributo (Dt 7:9).
\item[\emph{Haftará}] (\heb{הַפְטָרָה}). Lectura de los Profetas en la sinagoga, después de la Torá.
\item[\emph{Hekim}] (\heb{הֵקִים}). ``Levantar'', poner en pie: cumplir una palabra (1~R 8:20; Neh 9:8).
\item[\emph{Jésed}] (\heb{חֶסֶד}). Amor leal del pacto; pareja fija de \emph{emuná} (Sal 89).
\item[\emph{Ma'al}] (\heb{מַעַל}). Infidelidad, traición a lo sagrado (Nm 5:12).
\item[\emph{Milé}] (\heb{מִלֵּא}). ``Llenar'': la mano llena lo que la boca dijo (1~R 8:15, 24).
\item[\emph{Nafal}] (\heb{נָפַל}). Caer: ``no cayó una palabra'' (Jos 21:45; 1~R 8:56).
\item[\emph{Ne'emán}] (\heb{נֶאֱמָן}). Firme, seguro, fiel; participio nifal de \emph{amán}.
\item[\emph{Pístis}] (\grc{πίστις}). Fe, fidelidad (Ro 3:3; Gá 5:22).
\item[\emph{Pistós}] (\grc{πιστός}). Fiel, digno de confianza; también creyente (1~Co 1:9).
\item[\emph{Shoqed}] (\heb{שֹׁקֵד}). Vigilante: Dios ``vela'' sobre su palabra (Jer 1:12), con juego de palabras con \emph{shaqed}, almendro.
\end{description}

